\documentclass[11pt]{article}
\usepackage[margin=1in]{geometry}
\usepackage{amsmath,amssymb,amsthm}
\usepackage{hyperref}
\hypersetup{hypertexnames=false,hidelinks}
\usepackage{booktabs}
\usepackage{xcolor}
\usepackage{graphicx}
\usepackage{enumitem}
\usepackage{microtype}
\usepackage{array}
\usepackage{float}
\usepackage{longtable}
\newcolumntype{L}[1]{>{\raggedright\arraybackslash}p{#1}}

\title{When Does an LLM Forecasting Benchmark Measure Forecasting?\\
Source Currency, Label-Time Validity, and Equal-Information Controls}
\author{Daniel Alami}
\date{June 20, 2026}

\begin{document}
\maketitle

\begin{abstract}
LLM forecasting benchmarks often score model calls before establishing whether the scored row is a valid forecast. A row may ask a model to recover an answer already visible to its generation, use an outcome label from a later data vintage, or compare against a market or human baseline measured under a different information state. We define the forecast row as the unit of evidence and introduce three validity requirements: source currency, label-time validity, and equal-information baselines. The database contains more than 20{,}000 persisted calls, but the inferential units in this paper are not calls: the central source-currency and market-control tests use 24--80 contracts, pairwise ranking uses 24 non-tie pairs, and prompt-intervention checks use 90--120 contract-condition blocks. In a matched panel of 80 Manifold contracts, rows after the model cutoff are substantially harder than rows before the cutoff or visible in the model's sources (\(+0.191\) Brier in aggregate, paired-stratum delta \(+0.216\), permutation \(p=0.0004\), BH \(q=0.0031\), BY \(q=0.0115\)). The strict market comparisons remain small, so we use them to bound claims rather than to estimate a general market effect. On 24 Polymarket contracts, the four-family panel scores \(0.268\) Brier versus the market's \(0.073\) (\(p=0.0068\), BH \(q=0.0163\), BY \(q=0.0616\)). In a 24-contract Manifold slice, the market Brier is lower but the paired test is inconclusive; a separate 32-contract Manifold same-day freeze expansion scores \(0.215\) for the five-family panel versus \(0.136\) for the market (\(p=0.0048\), BH \(q=0.0163\), BY \(q=0.0616\)). Smaller Manifold horizon checks are reported only as overlapping sensitivity checks. These controls do not establish a general result about markets and models, and they do not support raw model panel superiority: two strict slices have much lower market Brier under the raw paired tests and BH correction, but the BY column treats them as sensitivity rather than arbitrary-dependence significance at \(0.05\); the third strict slice is inconclusive. After validity screening, model signal remains narrow. A selected rule that tempers very small model probabilities improves those estimates on eligible rows, and pairwise comparisons achieve \(0.750\) accuracy over 24 non-tie pairs when the pairs are balanced across sources. One expert-training prompt improves Brier in a completed 600-call Gemini experiment on public questions, but a 591/600-call Claude run remains underpowered and below the replication gate and a staged Codex+DeepSeek check does not reproduce the effect. We therefore treat it as a Gemini-specific candidate, not a general prompting method. The contribution is forecast row validity: a practical documentation layer for source, label, and comparator timing; an empirical demonstration that these checks change conclusions; a power-aware re-audit that reclassifies underpowered prior results instead of treating them as nulls; a consolidated multiplicity and effective-denominator audit; and a companion benchmark design that scores row validity, equal-information comparison, calibration, relative judgment, intervention, choosing among model families, open model replication, and public low-overlap replication as separate tracks.
\end{abstract}

\section{Introduction}

This paper introduces forecast row validity as a documentation standard for LLM forecasting benchmarks. The motivating questions are simple: when is a scored row actually measuring future-event prediction, and what conditions are needed before model outputs can be interpreted or used?

The starting point is an estimand problem. A forecasting leaderboard can assign a proper score to every row and still mix three different tasks: predicting a future event, recovering an answer already latent in the model's sources, and comparing against a human or market that saw different information. The central result is that LLM forecast output becomes interpretable only after the information state of each row is made explicit. A benchmark row is suitable for broad conclusions only after three conditions are documented. First, \emph{source currency}: the resolved answer was not already visible in the sources available to the model generation being tested. Second, \emph{label-time validity}: the outcome label matches what would have been knowable under the data vintage at forecast time. Third, \emph{equal-information baselines}: human or market comparison bars are measured on the same contract under the same information rule before outcome. In our data, the main Manifold source-currency audit is strong enough to report as a measurement result for that source. It is not yet a source-general prevalence estimate. Broad claims that the models beat humans or markets are not supported: the database has 170 typed external market baseline rows and 119 ingested equal-information market rows, but the completed strict controls are still modest. They include 24 Polymarket contracts, a 24-contract Manifold fill, a separate 32-contract Manifold same-day freeze expansion, and smaller Manifold one-day, two-day, and seven-day checks on overlapping rows. We use them as strict controls on the present evidence, not as population estimates of market or LLM performance. Within that scope, Polymarket has lower market Brier than the four-family model panel on the same 24 contracts, the first Manifold fill has lower market Brier but an inconclusive paired test, and the 32-contract Manifold expansion again has lower market Brier under the raw paired test and BH correction, with BY sensitivity reported below. The smaller one-day, two-day, and seven-day Manifold checks are sensitivity rows because they overlap the same source and remain too small for a population claim.

The empirical order matters. A black-box LLM forecaster first has to be tested on rows that are valid for its generation. Only then does it make sense to ask which emitted channel improves scoring for a particular family and source: point estimate, worry scalar, bid-ask spread, self-predicted Brier interval, reference-class base rate, or cross-family disagreement. Only after those two questions can we ask whether a prompt, abstention rule, review rule, or allocation rule improves Brier or utility against explicit controls.

The paper therefore has one main argument with two parts. The validity part is that source currency, label-time, and equal-information checks change what one is allowed to conclude from LLM forecasting benchmarks. The applied part is narrower: some model information survives the validity checks, but it appears through specific interfaces rather than a general forecasting ability. Very small model probabilities can be tempered on forward-looking rows that pass the source currency check; pairwise comparisons can rank harder and easier contracts better than chance when the pairs are balanced across sources; and one Gemini prompt improves Brier against bare and placebo prompts in a balanced comparison on public questions. That prompt result is not treated as a general method because the 591/600-call Claude run remains underpowered and below gate, and Codex+DeepSeek does not reproduce it. Differences across model families also create measurable headroom, but the paper treats that headroom as a target for future selection methods, not as a present decision rule. The practical output is a small set of controlled uses: a deterministic calibration rule for one row class, a ranking interface for relative judgment, a Gemini prompt result to replicate, and negative controls showing where not to spend inference budget.

This framing also explains why the paper is not organized as a single model leaderboard. The same score can mean different things depending on the row. A correction for very small probabilities is useful only after source and label checks; a pairwise comparison can be useful for ranking without being a calibrated probability; a prompt can improve one family while failing to replicate in others; and a market comparison is meaningful only when the market price is frozen at the same information time. The paper's contribution is to keep those cases separate and score each on the comparison it is actually allowed to answer.

\paragraph{Contributions.} The paper makes five contributions. First, it states a documentation test for LLM forecasting rows: a row is not broad forecast evidence unless source currency, label-time validity, and equal-information comparison status are known. Second, it gives an empirical source-currency audit in which Manifold rows after the model cutoff are substantially harder than matched rows before the cutoff or visible in the model's sources. Third, it reports equal-information market controls that do not support raw model panel superiority in the current evidence while remaining too small to estimate market performance in general. Fourth, it separates current usable results from design evidence. The current usable results are a selected calibration rule for eligible rows and pairwise ranking balanced by source. The completed Gemini expert-training prompt is reported as a model-specific intervention candidate and a replication target, not as a general prompt method. Additional structured outputs and differences across model families are design evidence for future interfaces rather than current applied results. These results are not claims that LLMs are superior to markets or humans. Fifth, it treats the power-aware re-audit as a methodological result and converts the missing-evidence map into a companion benchmark design: row validity fields, equal-information comparators, calibration, relative judgment, prompt intervention, choosing among model families, open model replication, and public low-overlap substitute tracks are evaluated separately rather than collapsed into one leaderboard.

\paragraph{Broader measurement route.} The paper also makes the broader measurement route concrete without claiming prevalence across the field. It does not report a failure rate across the field. The public ForecastBench audit scores 70 processed-forecast files over 521 resolved binary row keys and 230 event family keys, finds equal-information market slices in 68 files, and finds only 6 files beating the prior-day market baseline before and after event family capping. A second public ForecastBench audit on the 2024 human-comparator round scores 141 files over 7{,}259 row keys and 766 event family keys; the human-super and public aggregate files each have 577 resolved non-imputed rows, but each has only two strict equal-information market rows. A Prophet Arena source-access pilot fetches 68 public task rows across four sample releases, including 26 resolved rows, but the fetched samples contain no submitted forecast probabilities or same-time baseline probabilities; a public repository check over five AI Prophet repositories finds no committed Prophet Arena submission, leaderboard, or per-model trace archive. The PredictionMarketBench replay-row pilot reconstructs 370{,}254 same-time market-baseline rows but finds no stored model forecast rows in the released episodes; the PolyBench pilot verifies the repository/schema surface but cannot score the linked database from the noninteractive release path. These public audits show how the row validity machinery extends beyond this study, while leaving prevalence and conclusion-change rates as future claims.

\paragraph{Benchmark design implication.} The companion benchmark implied by the audit has eight separate tracks: row validity, equal-information comparators, calibration, relative judgment, intervention, choosing among model families or outside review, open model replication, and a public low-overlap substitute. Table~\ref{tab:benchmark-blueprint} gives the track structure. Keeping those tracks separate is part of the claim. A system can improve a calibrated probability without beating a market, rank pairs without supplying a reliable absolute probability, or show that model families make different errors without providing a working selection rule. During a new study, the same split tells the researcher which claim a packet can test before outcomes are scored: market comparison, calibration rule, relative-judgment interface, prompt intervention, selection rule, open model replication, or low-overlap diagnostic. The benchmark design is therefore both a way to collect the decisive missing evidence and a guard against turning one successful track into a broader claim. It is not a current claim about a measured failure rate across the field.

The row validity track is implemented as a required row schema, expanded in Table~\ref{tab:field-audit-protocol}: source currency, label-time validity, equal-information comparator timing, effective sample size, and decision-rule status. That schema is what makes the benchmark useful before scoring rather than only after a result looks surprising.

Used prospectively, the benchmark becomes an experiment-planning tool. Before calls are run or outcomes are known, each track asks what would make a positive result uninformative. If the answer is source visibility, label vintage, comparator timing, prompt length, family mix, or market overlap, that variable has to be recorded before scoring. This is the useful form of the inversion in practice: the question changes the design while there is still time to add the missing field or control, and a positive result is informative only when the simpler explanation it invites has already been made measurable.

\begin{table}[htbp]
\centering
\small
\begin{tabular}{L{0.22\linewidth} L{0.36\linewidth} L{0.32\linewidth}}
\toprule
Track & Question tested & Primary output \\
\midrule
Row validity core & Does each row state what the model could have known and which label is admissible? & Source-currency, label-time, and event family eligibility labels. \\
Equal-information comparators & Does the model beat a human, crowd, or market measured on the same contract at the same pre-outcome time? & Paired proper-score deltas with event family and source caps. \\
Probability calibration & Does a simple adjustment improve model probabilities on rows that pass the source currency check? & Calibrated versus raw Brier deltas by source, family, and cutoff relation. \\
Relative judgment & Are pairwise comparisons a better interface than direct probabilities? & Pairwise accuracy, utility, and predeclared graph-calibrated probabilities when justified. \\
Intervention & Does a prompt or procedure beat bare, placebo, calibrated bare, and matched comparator baselines? & Paired Brier deltas with source and family replication checks. \\
Choosing or reviewing model families & Can observable features or outside review recover the best family after cost? & Cost-adjusted Brier or utility against simple pools and calibrated baselines. \\
Open-model replication & Which scoped findings survive outside proprietary-provider snapshots? & Replication or bounded failure on public rows with open models. \\
Public low-overlap substitute & Do low-overlap elicitation diagnostics reproduce on a releasable corpus? & Channel and prompt diagnostics after novelty, source, topic, length, and horizon are separated. \\
\bottomrule
\end{tabular}
\caption{Companion benchmark tracks. The design keeps row validity, equal-information comparison, calibration, ranking, intervention, choosing among model families, open model replication, and low-overlap replication as separate scored tasks rather than a single aggregate leaderboard.}
\label{tab:benchmark-blueprint}
\end{table}

\paragraph{Positioning.} The closest related work falls into several groups: future-question benchmarks such as ForecastBench and Prophet Arena~\cite{karger2024forecastbench,yang2025prophetarena}, system papers such as AIA Forecaster~\cite{alur2025aiaforecaster}, belief-updating benchmarks such as EVOLVECAST~\cite{yuan2025evolvecast}, numerical forecast-interval benchmarks such as QuantSightBench~\cite{qin2026quantsightbench}, automated question-generation and resolution work~\cite{bosse2026automatingforecasting}, confidence-elicitation and fictional-market framing work~\cite{xiong2024confidence,todasco2025fakepredictionmarkets}, and market or replay environments such as Prediction Arena, PolyBench, PredictionMarketBench, Foresight Arena, MarketBench, and Reppo-style market infrastructure~\cite{zhang2026predictionarena,cheng2026polybench,arora2026predictionmarketbench,nechepurenko2026foresight,fradkin2026marketbench,reppo2026}. These systems evaluate broader capabilities: generating forecasts, updating after new information, expressing calibrated intervals over continuous quantities, generating and resolving questions, eliciting confidence or wagers, trading with real or replayed market frictions, or coordinating through market-like designs. Their scores can mix probability accuracy with retrieval timing, execution quality, liquidity, fees, position sizing, interval width, confidence reporting, automated resolution quality, and the moment at which a human or market comparison was sampled. Nearby work also studies market relationships and relative judgment: Semantic Trading uses agentic AI to cluster prediction markets and identify correlated or anti-correlated market pairs~\cite{capponi2025semantictrading}, while a fully prospective venture tournament finds strong model rankings from pairwise comparisons against human managers and investors~\cite{csaszar2026strategicforesight}. These results are closest to this paper's pairwise-ranking result, while our main contribution remains different. This paper studies the narrower evidentiary unit underneath those comparisons. Before a system-level score is interpreted as forecasting evidence, each row has to state what the model could have known, which label vintage is admissible, and whether the comparator was measured under the same information rule. The paper is therefore closest in spirit to recent evaluation-warning work on temporal leakage and benchmark extrapolation~\cite{paleka2025pitfalls}, but adds a scored empirical audit and controlled tests for using model output after the validity checks.

\paragraph{Evidence discipline.} Table~\ref{tab:claim-map} states the paper's main claims and their limits. It is included near the front because the main risk in this area is not a missing benchmark but a wrong comparison: rows that are source-visible, labels that use a later data vintage, market bars measured at a different information time, or decision rules adopted from diagnostic correlations. Read in the opposite direction, the test is simple: if the answer, label, or comparator comes from a different information state than the forecast, the row may still be useful for diagnosis, but it cannot support a broad forecasting comparison.

\begin{table}[htbp]
\centering
\small
\begin{tabular}{L{3.0cm} L{4.4cm} L{5.5cm}}
\toprule
Claim unit & Evidence in this paper & Conservative interpretation \\
\midrule
Validity checks & Source-currency, label-time, and equal-information audits over the database & A row without these checks is not broad forecast evidence. \\
Raw LLM panels vs markets & Three strict equal-information controls: Polymarket and the 32-contract Manifold expansion have much lower market Brier under raw paired tests and BH correction, with BY sensitivity just above \(0.05\); the 24-contract Manifold fill is inconclusive. & Does not support raw model panel superiority in the current evidence; too small for a general market-performance estimate. \\
Calibration for very small probabilities & Rule for very small probabilities improves every family on the public-domain panel and improves the forward-looking slice & Probability use for rows that pass the source currency check, not retrospective correction. \\
Pairwise ranking & Pairwise ranking balanced by source and partial probability-translation evidence & Ranking/tournament support; absolute-probability use needs larger controls. \\
Expert-training prompt & 600-call Gemini public question comparison: expert-training improves paired Brier versus bare, length-matched placebo, and same row calibrated bare forecasts; two other structured prompt variants do not beat placebo & A Gemini-specific result and replication target; current market overlap has lower market Brier, the 591/600-call Claude run remains underpowered and below the replication gate, and a 448-call Codex+DeepSeek staged run does not clear the bare, placebo, sign-test, or source-split checks. \\
Evidence fields and family choice & Mixed structured-field tests; best-family-in-hindsight headroom not recovered by simple allocation rules & Headroom is a design target, not a predictive asset until an observable selector recovers it prospectively. \\
Generic prompt intervention & Selective action, self-revision, and diagnostic allocation mostly fail controls & Evidence against unvalidated prompt-only interventions. \\
\bottomrule
\end{tabular}
\caption{Main claims and limits. The central contribution is the validity layer plus the controlled tests for using LLM forecast information once invalid comparisons are removed.}
\label{tab:claim-map}
\end{table}

\begin{figure}[htbp]
\centering
\small
\begin{tabular}{@{}L{0.22\linewidth} L{0.34\linewidth} L{0.34\linewidth}@{}}
\toprule
\textbf{Stage} & \textbf{Question answered} & \textbf{Consequence} \\
\midrule
Forecast row & What probability, source, timestamp, and side channels did the model emit? & Defines the unit to be scored. \\
Validity checks & Was the answer source-visible, and is the label valid at the relevant data vintage? & Decides whether the row can count as forecast evidence. \\
Equal-information baseline & What did a market or human comparator know at the same pre-outcome time? & Decides whether raw model probabilities add value under the same information rule. \\
Controlled use & Which model-derived output improves a proper score or ranking task after the checks above? & Admits only calibration on rows that pass the source currency check, pairwise ranking, or future differences across model families rules that beat controls. \\
\bottomrule
\end{tabular}
\caption{Evidence flow. The paper's unit is a forecast row with documented validity checks, an equal-information baseline, and a controlled test of which model-derived output improves scoring or ranking.}
\label{fig:evidence-flow}
\end{figure}

The paper is empirical first and theoretical second. We ran a calibration study across five model families and two corpus classes, with more than 30 measured findings and pre-registered tests. The main text keeps the forecasting argument in front: validity checks, equal-information baselines, and controlled uses of model output. Bias-transfer and prompt-stability results are treated as secondary diagnostics because they help explain why generic prompting is unreliable, but they do not carry the paper's main claim.

\paragraph{Statistical audit.} Every reported test is checked with the same power-aware procedure: Fisher-$z$ sample-size computation before an experiment runs, three possible outcomes after scoring (supported, ruled out at the target effect size, or underpowered), equivalence testing instead of treating ``$p > 0.05$'' as evidence of no effect, and leave-one-out $R^2$ at small $N$. The manuscript reports one global BH-FDR correction over all p-values used for its main statistical comparisons and a Benjamini-Yekutieli column as a conservative robustness check for arbitrary dependence across overlapping panels. Table~\ref{tab:multiplicity-audit} gives the effective denominator, raw p-value, BH \(q\), and BY \(q\) for each test. For the Polymarket and 32-contract Manifold market rows, BH gives \(q=0.0163\) while BY gives \(q=0.0616\); we therefore treat those rows as strong current controls and BY sensitivity evidence, not as arbitrary-dependence-significant market effects at \(0.05\). The class column separates diagnostic, exploratory, replication, and continuation rows; none of the exploratory or sensitivity rows is promoted solely by a multiplicity-adjusted threshold. Under the broader audit, $8$ of $12$ prior nulls in the study became underpowered rather than negative findings, and $5$ cross-corpus claims narrowed to corpus-specific results. That procedure is a methodological contribution because the same overcalling problem plausibly affects existing $N{=}5$--$20$ cross-corpus headlines. A second-order application is corpus-validity drift: a benchmark whose resolution dates were post-cutoff for publication-era LLMs can become pre-cutoff for the next model generation. We apply this point to the Halawi 2024 dataset~\cite{halawi2024approaching} in $\S$\ref{sec:reaudit}.

\begingroup
\scriptsize
\setlength{\tabcolsep}{2pt}
\renewcommand{\arraystretch}{0.9}
\begin{longtable}{@{}L{0.19\linewidth} L{0.09\linewidth} L{0.14\linewidth} L{0.09\linewidth} L{0.15\linewidth} L{0.06\linewidth} L{0.06\linewidth} L{0.06\linewidth}@{}}
\caption{Consolidated multiplicity and effective-denominator audit. Lower Brier is better. Positive panel-minus-market or expert-minus-market values mean the model or prompt is worse than the market. BH values use one global Benjamini-Hochberg family over the 24 p-values reported here; BY values apply Benjamini-Yekutieli as an arbitrary-dependence robustness check.}
\label{tab:multiplicity-audit}\\
\toprule
Test & Class & Unit & Effective \(N\) & Effect & \(p\) & BH \(q\) & BY \(q\) \\
\midrule
\endfirsthead
\toprule
Test & Class & Unit & Effective \(N\) & Effect & \(p\) & BH \(q\) & BY \(q\) \\
\midrule
\endhead
\bottomrule
\endfoot
Manifold source-currency audit & Diag. & Contract / paired stratum & 80 / 15 & Post-minus-pre Brier \(+0.191\); paired delta \(+0.216\) & 0.0004 & 0.0031 & 0.0115 \\
Polymarket market control & Diag. & Contract & 24 & Panel-minus-market \(+0.195\) & 0.0068 & 0.0163 & 0.0616 \\
Manifold market control & Diag. & Contract & 24 & Panel-minus-market \(+0.0377\) & 0.5431 & 0.6207 & 1.0000 \\
Manifold same-day freeze & Diag. & Contract & 32 & Panel-minus-market \(+0.0787\) & 0.0048 & 0.0163 & 0.0616 \\
Manifold one-day freeze & Sens. & Contract & 18 & Panel-minus-market \(+0.1026\) & 0.0122 & 0.0244 & 0.0921 \\
Manifold two-day freeze & Sens. & Contract & 10 & Panel-minus-market \(+0.1225\) & 0.0152 & 0.0281 & 0.1060 \\
Manifold seven-day freeze & Sens. & Contract & 7 & Panel-minus-market \(+0.0346\) & 0.5045 & 0.6054 & 1.0000 \\
Very-small-probability rule vs raw panel & Expl. & Contract panel & 132 & Raw-minus-tempered Brier \(+0.0296\) & 0.0062 & 0.0163 & 0.0616 \\
Very-small-probability rule, forward-looking rows & Diag. & Calls / contracts & 120 / 40 & Tempered-minus-raw Brier \(-0.0253\) & 0.0688 & 0.1032 & 0.3897 \\
Very-small-probability rule, source-visible rows & Diag. & Calls / contracts & 120 / 40 & Tempered-minus-raw Brier \(+0.0350\) & 0.0002 & 0.0024 & 0.0091 \\
Pairwise ranking vs random & Expl. & Unique non-tie pair & 24 & Accuracy \(0.750\); utility \(+0.583\) & 0.0044 & 0.0163 & 0.0616 \\
Pairwise ranking vs source control & Expl. & Unique non-tie pair & 24 & Utility \(+1.583\) & 0.0002 & 0.0024 & 0.0091 \\
Pairwise probabilities vs raw panel & Cont. & Contract & 24 & Candidate-minus-baseline Brier \(-0.0061\) & 0.7351 & 0.7436 & 1.0000 \\
Pairwise probabilities vs low-probability rule & Cont. & Contract & 24 & Candidate-minus-baseline Brier \(-0.0158\) & 0.3229 & 0.4559 & 1.0000 \\
Pairwise probabilities vs market & Cont. & Contract & 24 & Candidate-minus-market Brier \(-0.0130\) & 0.5783 & 0.6309 & 1.0000 \\
Gemini expert-training vs bare & Expl. & Contract-condition block & 120 & Expert-minus-bare Brier \(-0.0606\) & 0.0005 & 0.0031 & 0.0115 \\
Gemini expert-training vs placebo & Expl. & Contract-condition block & 120 & Expert-minus-placebo Brier \(-0.0243\) & 0.0067 & 0.0163 & 0.0616 \\
Gemini expert-training vs tempered bare & Expl. & Contract-condition block & 120 & Expert-minus-tempered-bare Brier \(-0.0525\) & 0.0022 & 0.0103 & 0.0390 \\
Gemini expert-training vs all matched markets & Diag. & Contract with market & 51 & Expert-minus-market Brier \(+0.1510\) & 0.0110 & 0.0239 & 0.0904 \\
Gemini expert-training vs equal-information markets & Diag. & Contract with market & 33 & Expert-minus-market Brier \(+0.0931\) & 0.0351 & 0.0601 & 0.2271 \\
Claude replication vs bare & Repl. & Paired block & 115 & Expert-minus-bare Brier \(-0.0036\) & 0.7436 & 0.7436 & 1.0000 \\
Claude replication vs placebo & Repl. & Paired block & 112 & Expert-minus-placebo Brier \(-0.0042\) & 0.0530 & 0.0848 & 0.3202 \\
Codex+DeepSeek replication vs bare & Repl. & Paired block & 90 & Expert-minus-bare Brier \(+0.0072\) & 0.4704 & 0.5942 & 1.0000 \\
Codex+DeepSeek replication vs placebo & Repl. & Paired block & 90 & Expert-minus-placebo Brier \(+0.0652\) & 0.3891 & 0.5187 & 1.0000 \\
\end{longtable}
\endgroup

The generated audit report is mirrored in the root project:
{\scriptsize\path{llm-forecast-calibration-cross-corpus/evidence/reports/multiple_testing_effective_n_audit.md}.}

\paragraph{Reading order.} Section $\S$\ref{sec:setup} defines the panel, corpora, and verdict rules; $\S$\ref{sec:core-results} states the paper's core empirical results before the diagnostics. Sections $\S$\ref{sec:theoretical-frame}--$\S$\ref{sec:conditional} give the elicitation-surface and family-heterogeneity measurements that motivated the validity audits. Section $\S$\ref{sec:universal} reports the applied calibration and allocation consequences, including the calibration for very small probabilities and pairwise-ranking results. Section $\S$\ref{sec:controlled-use} synthesizes the controlled-use result under the same source, label-time, and market-baseline constraints. Section $\S$\ref{sec:bias-transfer} gives the secondary bias-transfer and prompt-stability diagnostics. The main measurement audit is collected in $\S$\ref{sec:reaudit}: corpus-validity drift, source currency checks, label-time checks, market-prior repair, and equal-information market baselines. The paper closes with explicit limits, next tests, a conclusion, and a reproducibility plan.

\section{Setup}
\label{sec:setup}

\subsection{Five LLM families}
Claude Opus 4.7 (Anthropic, via Claude Code CLI), GPT-5.5 and GPT-5.4-mini (OpenAI, via Codex CLI), Gemini 2.5 Flash (Google GenAI API), DeepSeek Chat (DeepSeek API). All API calls $\text{temperature}{=}0$ unless explicitly noted. All paths were audited for web-search-tool access; all prompts included an explicit instruction not to use web search.

\subsection{Two corpora, four confounded axes}
\label{sec:corpora}

We use a \emph{low-overlap corpus} of $N{=}15$ bespoke questions drawn from a private research domain with effectively zero LLM-pretraining overlap (formal-proof tactic combinations, state-transition questions from a private research workflow, and ground-truth-confirmed forecasting contracts), and a \emph{public-domain corpus} of $N{=}42$ contracts from prediction markets (Manifold, Polymarket), stock-close thresholds, and ETF moves.

\begin{table}[htbp]
\centering
\small
\begin{tabular}{l L{5.5cm} L{5.5cm}}
\toprule
Axis & Low-overlap ($N{=}15$) & Public-domain ($N{=}42$) \\
\midrule
LLM pre-training overlap & Effectively zero & Substantial \\
Question length & 247 chars (35 words) & 99 chars (17 words) \\
Question structure & Multi-conditional, code/math symbols, parenthetical embedding & Atomic, simple syntactic structure \\
Implicit base rate & Ill-defined (bespoke operations) & Public-market priors well-trained \\
\bottomrule
\end{tabular}
\caption{Corpus contrast. The low-overlap and public-domain corpora differ on several axes at once, so cross-corpus differences are treated as diagnostic rather than as a single-factor causal estimate.}
\label{tab:corpus-contrast}
\end{table}

The four axes are confounded in the current data. Low-overlap-corpus questions are long, complex, unfamiliar, and without public priors \emph{together}; public-domain questions are short, simple, familiar, and public-prior-rich \emph{together}. The cross-corpus result ($\S\ref{sec:frequency-framing}$) is therefore evidence that \emph{at least one} of these axes matters; a 4-cell de-confounded corpus design that breaks the confound is the highest-priority methodological follow-up ($\S\ref{sec:limits}$).

The low-overlap corpus is private and not publicly releasable in raw form; the reproducibility plan ($\S\ref{sec:repro}$) addresses this through a sanitized release with neutral identifiers and a parallel public-but-niche-academic corpus matching the four-axis profile. It is not required for the paper's central source currency, label-time, equal-information, market control, or calibration claims on rows that pass the source currency check. Those claims are scored from the public database, market-history files, official-data label-time checks, and audit scripts listed in $\S\ref{sec:repro}$. Low-overlap rows are retained as secondary diagnostics about elicitation channels, prompt stability, and external generality.

\subsection{Forecast row validity and the comparison estimand}
\label{sec:row-estimand}

The paper's comparison unit is a scored forecast row, not a model call by itself. A row contains a question or contract, a model family, a forecast timestamp, the model's cutoff or retrieval window, a probability, an outcome label with data vintage or settlement rule, source metadata, and any human or market comparator with its own timestamp. A broad model-vs-comparator claim is estimated only on rows for which the validity indicator \(V_{\mathrm{row}}\) is one: the resolved answer was not source-visible to the model at generation time; the outcome label is admissible under the relevant data vintage or settlement rule; and the comparator probability, if used, was measured on the same contract under the same pre-outcome information rule. Rows with \(V_{\mathrm{row}}=0\) are still useful for diagnosing retrieval, source familiarity, label revision, or market-timing effects, but they are not in the denominator for broad claims that an LLM, human, crowd, or market is the better forecaster. This is the estimand behind the tables below: differences in Brier score are interpreted only after the row's information state is fixed.

The same rule is applied before a new packet is scored. For each planned result, we write the plain counter-explanation that would make a positive result uninformative: a comparator observed at the wrong time, an answer visible in the source, a later label vintage, a prompt-length or placebo effect, an imbalanced pair orientation, or a source mix that makes the denominator misleading. The experiment then has to include the row field or control that could rule out that explanation. The rule is meant to change the test while it is still being planned, not to explain a result after the score is known. If it cannot be satisfied, the result may still be useful for diagnosis or design, but it is not counted as a broad performance claim.

\subsection{Pre-registration and power-aware verdicts}

Every experiment is pre-registered before launch with hypothesis, null, target effect size, $n_{\text{required}}$ (computed via Fisher-$z$ at $\alpha{=}0.05$ two-tailed, $80\%$ power, Spearman correction $+6\%$), falsifiers, and success criterion. Resolution uses three possible outcomes:

\begin{itemize}[leftmargin=*]
\item \textbf{Supported}: 95\% CI on observed $\rho$ excludes 0.
\item \textbf{Ruled out at the target effect size}: 95\% CI wholly within $(-\text{target}, +\text{target})$; data rules out a meaningful effect.
\item \textbf{Underpowered}: observed $|\rho|$ below detectability at the run $N$; CI wide.
\end{itemize}

Brier-delta tests use paired-contract sign-flip permutation with a 90\% bootstrap CI plus a BIC-approximation Bayes factor. ``No effect'' claims also require TOST equivalence at a pre-stated bound. Power thresholds at $\alpha{=}0.05$ / power$=0.80$: $N{=}15 \Rightarrow |\rho| \geq 0.68$; $N{=}42 \Rightarrow 0.43$; $N{=}91 \Rightarrow 0.30$.

\section{Core empirical results}
\label{sec:core-results}

The rest of the manuscript gives the diagnostics behind these claims. The core evidence is the following.

\begin{table}[htbp]
\centering
\scriptsize
\begin{tabular}{@{}L{0.17\linewidth} L{0.20\linewidth} L{0.39\linewidth} L{0.16\linewidth}@{}}
\toprule
Result & Status and denominator & Evidence & Paper use \\
\midrule
Source-currency validity matters & Diagnostic; 80 contracts / 15 paired strata. & Matched Manifold panel: 80 contracts / 240 tool-free calls; post-minus-pre Brier \(+0.191098\), paired-stratum delta \(+0.2155\), \(p=0.0004\), BH \(q=0.0031\), BY \(q=0.0115\). & Manifold measurement result; source-general prevalence remains unmeasured. \\
Current market controls do not support raw panel superiority & Diagnostic controls; 24 Polymarket contracts, 24 Manifold contracts, and one 32-contract Manifold expansion. Smaller horizon fills are sensitivity rows only. & Polymarket: panel \(0.267758\) vs market \(0.072964\), \(p=0.0068\), BH \(q=0.0163\), BY \(q=0.0616\). Manifold fill: market Brier lower but paired test inconclusive (\(p=0.5431\)). Manifold same-day expansion: panel \(0.214665\) vs market \(0.135951\), \(p=0.0048\), BH \(q=0.0163\), BY \(q=0.0616\). & Boundary: current controls do not support claims that LLMs are superior to markets or humans. \\
Calibration for very small probabilities improves point probabilities under scope & Exploratory selected rule; 132 contract panels. & Raw mean-panel remains \(+0.029598\) Brier worse than the correction for very small probabilities on rows with adequate source and label-time documentation (\(p=0.0062\), BH \(q=0.0163\), BY \(q=0.0616\)), but the rule regresses source-visible rows. & Selected calibration rule for rows that pass the source currency check, not a universal correction. \\
Pairwise ranking gives controlled relative-judgment evidence & Exploratory; 24 unique non-tie pairs. & Pairwise comparison balanced by source: accuracy \(0.750\), utility \(+0.583\), \(p=0.0044\) vs random, BH \(q=0.0163\), BY \(q=0.0616\); source-control comparison has BY \(q=0.0091\). & Ranking support, not standalone probability. \\
Expert-training prompt passes a narrow public question test & Exploratory Gemini result; 120 contract-condition blocks; replications are continuation rows. & Completed 600-call Gemini comparison: expert-training improves paired Brier versus bare prompt (BH \(q=0.0031\), BY \(q=0.0115\)), length-matched placebo (BH \(q=0.0163\), BY \(q=0.0616\)), and calibrated bare forecast; audit-informed and failure-mode-specific prompts do not beat placebo. & Gemini-specific candidate; current market overlap has lower market Brier, the 591/600-call Claude run remains underpowered and below the replication gate, and Codex+DeepSeek does not reproduce the effect. \\
Generic prompt intervention mostly fails & Negative guidance; denominators vary by packet. & Selective action, generic self-revision, and diagnostic allocation fail or regress against controls. & Generic prompting is not enough. \\
\bottomrule
\end{tabular}
\caption{Core empirical results. The diagnostic sections explain how these results were found and where they fail; this table states the claims that remain after source, label-time, and market-baseline controls.}
\label{tab:core-results}
\end{table}

This ordering is deliberate. The paper's positive claims are downstream of the equal-information market controls. A controlled-use result is included only when it either improves a proper score on valid rows, improves a balanced comparison against bare and placebo prompts, supports a controlled ranking use, or identifies structure that a future decision rule could try to recover. Mechanisms that only change rationales, increase stated effort, or improve a source-visible retrospective slice are reported as diagnostics or failed interventions.

\begin{figure}[htbp]
\centering
\includegraphics[width=0.82\linewidth]{evidence/figures/equal_information_market_controls.png}
\caption{Equal-information market controls. Lower Brier is better. The plot is generated from the stored Polymarket replacement and Manifold history score reports by the figure script listed in the reproducibility section. Polymarket has much lower market Brier than the model panel; in the first Manifold fill the market has lower Brier, but the paired test is inconclusive. These controls do not support a raw model panel superiority reading for the current evidence, but do not estimate a general market control effect.}
\label{fig:equal-info-bars}
\end{figure}

\section{Why emitted forecast channels differ by family}
\label{sec:theoretical-frame}

Before reporting the elicitation findings, we name the structural distinctions that organized the measurements. These were not stated in advance; they are a compact reading of the body of findings the study produced. Their role in this paper is explanatory: they help explain why side channels and prompt interventions are family- and source-conditional, not evidence that LLMs beat markets.

\paragraph{Elicitation surface.} An LLM forecaster, when asked, emits a set of channels: at minimum a point estimate $p_{\text{success}}$; optionally a tail-worry scalar (``how worried are you?''), a bid-ask spread $(p_{\text{buy yes max}}, p_{\text{sell yes min}})$, a self-predicted Brier interval $(b_{\text{lo}}, b_{\text{hi}})$, an asymmetric-loss decision threshold under specified false-positive vs. false-negative cost regimes, or an outside-view base rate from analogous prior contracts. Training-mix differences across model families produce different elicitation surfaces, even when prompted identically. $\S\ref{sec:bid-ask-spread}$ introduces the bid-ask spread as a diagnostic error-warning channel that is elicitable on the low-overlap corpus; $\S\ref{sec:channel-orthogonality}$ shows the three uncertainty channels are statistically independent on a paired panel, motivating per-family channel analysis; $\S\ref{sec:multi-channel-r2}$ shows only one family generalises a multi-channel decomposition under leave-one-out cross-validation.

\paragraph{Behavioral structured-field interpretation.} We do not observe the hidden activations of the closed models used here, so our channel claims are behavioral rather than mechanistic-interpretability claims. The measurable object is the forecast row and its scored side fields, not the generated rationale by itself. Recent latent-prediction theory shows that learning or predicting hidden representations can avoid the sample-complexity cost of token-level prediction on hierarchical data~\cite{korchinski2026latents}. Our setting is not training-time latent prediction, but the comparison clarifies the measurement problem: uncertainty channels and source-bound evidence fields can be scored against outcomes, while free-form rationale text is a noisy surface output. In two small follow-up tests across two families, structured evidence fields beat free prose on mean Brier, while the stricter two-step variant did not consistently beat a same-turn field. The combined means were baseline $0.171038$, free prose $0.146103$, two-step field $0.103278$, and same-turn field $0.098425$. A later placebo-control test weakened the claim: among the 30 rows with complete structured outputs, baseline mean Brier was $0.078000$, two-call prose $0.107254$, same-turn field $0.110300$, free prose $0.122767$, and two-step field $0.149921$; ten additional Codex rows failed at runtime before emitting forecasts. The current data justify a larger paired study of structured evidence fields, but not their use as a scoring rule.

\paragraph{Bias-transfer diagnostic.} We also tested whether human-style cognitive-bias categories predicted LLM forecast distortions. The useful distinction is representational: some biases have a large natural-language footprint, some depend on human utility or affect, and some are heavily discussed in case-study text even though the model has no corresponding utility function. This distinction helped organize prompt-risk and family-sensitivity tests, but it did not establish a forecasting intervention. In the 180-call anti-bias-prompt test, the average collapse pattern was underpowered, the class-label shuffle was null ($p=0.5387$), and the raw-gap-adjusted coefficient for text-discussed motivational cases was negative ($-0.076587$, $p=0.0025$). A matched audit using the stored rows also found poor support overlap: at raw-gap caliper $0.05$, within-family matching left only 16 with-replacement / 15 no-replacement pairs and flipped the estimated collapse effect negative ($-0.072750$, $p=0.0008$; greedy no-replacement $-0.077561$, $p=0.0006$), while the no-caliper positive estimate relied on large raw-gap distances. The supported claim is therefore diagnostic: bias-category measurements help explain prompt fragility, but a causal cognitive-debiasing intervention requires new matched strata or randomization.

\paragraph{Family/post-training overlay.} Family identity and post-training choices reshape per-family channel surfaces, but this paper observes only cross-family contrasts and does not identify the causal role of RLHF. Family-level worry/error sign differences ($\S\ref{sec:worry-direction-split}$), per-family channel-decomposition $R^2$ differences ($\S\ref{sec:multi-channel-r2}$), and the self-assessed channel-choice failure where three of five families pick the worse channel at $p \leq 0.005$ each ($\S\ref{sec:conditional}$) all live on this axis. The empirical signature is that family identity, conditional on channel and contract, is itself a central predictor.

\paragraph{Scoring discipline across diagnostics.} Every diagnostic above is evaluated with the same power-aware procedure: Fisher-$z$ power calculation before launch, three possible outcomes after scoring, equivalence testing, global BH-FDR with BY robustness where many reported tests are compared, and LOO-$R^2$ at small $N$. This discipline produced the study's most consequential retraction: 8 of 12 prior ``no effect'' claims became underpowered rather than negative findings, and the corpus-validity-drift filter (resolve\_date $>$ \texttt{max(panel\_cutoff)}) empties the most-cited 2024 LLM-forecasting benchmark for the current LLM generation ($\S\ref{sec:reaudit}$).


\section{Elicitation diagnostics: channels are conditional, not decision rules}
\label{sec:worry-foundation}
\label{sec:worry-direction-split}
\label{sec:cognitive-decoupling}
\label{sec:bid-ask-spread}
\label{sec:channel-orthogonality}
\label{sec:frequency-framing}
\label{sec:worry-mechanism}
\label{sec:multi-channel-r2}
\label{sec:conditional}

The uncertainty-channel experiments explain why raw probabilities alone are insufficient, but they are not themselves decision rules. We elicited worry scalars, bid-ask spreads, trajectory variance, frequency framings, self-predicted Brier intervals, outside-view base rates, and self-assessments about which channel the model expected to predict error. The pattern is consistent across these probes: LLMs emit useful side information about contracts, families, and sources, but the side information is conditional. It changes sign by family, weakens under task-difficulty controls, and often fails when moved from a low-overlap corpus to public-domain market questions.

\begin{table}[htbp]
\centering
\small
\begin{tabular}{L{0.23\linewidth} L{0.38\linewidth} L{0.31\linewidth}}
\toprule
Diagnostic & Main evidence & Manuscript role \\
\midrule
Worry scalar & Pool-level worry is a positive tail-risk signal at $N{=}590$, but per-family worry-Brier signs split. Claude and GPT-5.5 can be most worried when most right, while GPT-5.4-mini is direction-sensible. Topic-trigger and mean-regression probes explain the sign flip better than rationale length, wallclock time, or hedging vocabulary. & Use worry as a behavioral channel only after conditioning on family, corpus, and source. It is not a universal calibration transform. \\
Cognitive text versus calibration & Failure-mode words in rationales are weakly positive for the original trio but near zero for Gemini and DeepSeek at roughly $N{=}60$ per family. Stake framing raises worry for every family while leaving point probabilities nearly unchanged. & The model can produce diagnostic language without recomputing the forecast from that language. This supports interface discipline, not free-form chain-of-thought reliance. \\
Bid-ask and frequency probes & Bid-ask spread is positive on the low-overlap corpus for GPT-5.4-mini and DeepSeek, but public-domain replication is underpowered or sign-flipped. Frequency framing improves GPT-5.4-mini on low-overlap questions and does not transfer to the public question. & These are promising elicitation probes. They do not justify use across sources. \\
Orthogonality and multi-channel fits & Worry, spread, and trajectory variance are nearly uncorrelated at expression level. Early same-forecast $R^2$ survives leave-one-out only for GPT-5.5 on the $N{=}42$ public question; the larger five-family audit gives negative leave-one-out $R^2$ for channel-only Brier prediction across all five families. & Channels are distinct measurements, but distinct does not mean ready for use. They are useful for diagnosis and design of later scoring rules. \\
Conditional structure and family allocation & Sub-source decomposition flips channel signs across yfinance, Manifold, and Polymarket cells. Pooled channel value collapses after task-difficulty controls. Easy-contract cells show headroom, hard cells overfit, and naive mean/median panels do not beat the best single family. & Model information appears in family-by-source-by-contract interactions. Current observable allocation rules find headroom but do not yet recover it under controls. \\
\bottomrule
\end{tabular}
\caption{Uncertainty-channel experiments as diagnostics. The shared lesson is conditionality: emitted channels can locate signal, but they have not become probability rules across sources.}
\label{tab:elicitation-diagnostics}
\end{table}

The channel material has a narrow role in the argument. The diagnostics support one claim needed for controlled use: raw LLM forecasts contain more structure than a single probability, and that structure requires label-time-valid scoring, source and family splits, and comparison against simple baselines before use.

Several negative results are central. Channel-only decision rules that look attractive in small or low-overlap cells fail under larger leave-one-out checks or source controls. Models can identify contracts that feel leaky or risky without translating that feeling into better probabilities. Self-assessed channel selection can even be anti-coherent: in the low-overlap channel-choice probe, three of five families often selected the channel that was worst for their own error. These failures motivate the paper's emphasis on calibration on rows that pass the source currency check, pairwise ranking, source-specific checks, and future allocation tests rather than generic self-monitoring or prompt-only intervention.

\section{Additional all-channel diagnostics}
\label{sec:all-channel-diagnostics}

The all-channel prompts collected several side measurements in a single forecast pass or as paired conditions. Five additional findings emerge at $N{=}131$--$142$ public-domain:

\textbf{Rollback hurts.} A self-counterfactual prompt (``imagine you were trained only up to 2023, emit a rollback forecast'') worsens Brier for 4 of 5 families at $N{=}142$ paired ($\Delta$ Brier rollback$-$normal: claude $+0.064$, codex-5.5 $+0.034$, codex-5.4-mini $+0.028$, deepseek $+0.028$, gemini $-0.005$). Post-2023 training is net-useful even though families cannot introspect which contracts benefit. \emph{Practical caution:} avoid ``imagine you knew less'' framings on forecasting tasks without a heldout correction check.

\textbf{Leakage agreement without calibration.} On the same rollback prompt, all 10 pairwise $\rho(|\text{leak}_A|, |\text{leak}_B|)$ between families are $\geq +0.43$ (claude $\leftrightarrow$ codex-5.5 $\rho{=}+0.80$). Families agree on which contracts feel high-leakage. Per-family $\rho(|\text{leak}|, \text{err}^2)$ is nevertheless equivalent to zero for all five at the $\pm 0.30$ target. Self-reported leakage is a contract-structural signal, not a per-family calibration channel. It is useful as a corpus-curation filter for low-leakage LLM-fresh tests.

\textbf{DeepSeek's emitted decision threshold is calibrated.} In an asymmetric-loss decision-threshold prompt (commit-YES threshold under symmetric / YES-bad-3$\times$ / NO-bad-3$\times$ loss), Codex-5.4-mini emits the most Bayes-coherent thresholds (mean YES-bad $0.750$, exact match). For DeepSeek, $\rho(\text{threshold\_yes\_bad}, \text{err}^2) = -0.19$, $95\%$ CI $[-0.35, -0.03]$ at $N{=}142$. The emitted threshold tracks accuracy for one family. This is a novel diagnostic channel; applied use requires a heldout utility test.

\textbf{Stake-framing amplifies worry universally, leaves $p$ unchanged.} Same contract framed as ``\$1 stake'' vs ``\$100K stake'' across 103--142 paired contracts per family. Worry-scalar amplifies for every family ($\Delta\text{worry}$ from $+10.8$ deepseek to $+33.5$ claude on the 1--100 scale). But $|\Delta p_{\mathrm{success}}| < 0.03$ for every family. This cognitive-shift-without-calibration-shift result connects the cognitive-decoupling result to the stake-framing and risk-amplitude literature in human decision-making~\cite{holt2002risk}.

\textbf{Self-assessed channel choice fails for three of five families.} A channel-choice prompt (``which of your own uncertainty channels best predicts your error?'') reveals systematic anti-coherence at $N{=}15$ on the low-overlap corpus: claude, codex-5.5, and deepseek pick the worse channel $87\%$ of the time (binomial $p{=}0.005$ each). Within those picks, $\rho(\text{picked}, \text{err}^2)$ is negative while $\rho(\text{unpicked}, \text{err}^2)$ is strongly positive (claude $+0.51$, codex-5.5 $+0.53$, deepseek $+0.65$). A testable inversion pattern is: when these three families claim ``worry is my best channel,'' test spread instead (or the analogous flip). Codex-5.4-mini is coherent ($73\%$ right); gemini is random. Companion finding to the per-family worry sign-flip, on a new axis. This result is low-overlap-only at $N{=}15$; cross-corpus replication at $N{=}42$ external is pre-registered. The pattern (LLMs have contract-level structural intuitions without calibration to truth) recurs across the rollback, leakage, decision-threshold, and channel-choice results, the defining negative finding about LLM self-monitoring.

\section{Applied calibration and allocation consequences}
\label{sec:universal}

The channel and family diagnostics become useful only if they change a scored forecast or a decision rule. This section separates three evidence grades: robust enough for a calibration on rows that pass the source currency check view, useful as a design clue, and unsupported as an applied rule.

\begin{table}[htbp]
\centering
\small
\begin{tabular}{L{0.25\linewidth} L{0.43\linewidth} L{0.24\linewidth}}
\toprule
Pattern & Evidence & Status \\
\midrule
Overconfidence on very small probabilities & In the lowest forecast quintile, every family underpredicts YES outcomes; gaps range from \(+0.34\) to \(+0.72\). & Basis for the rule for very small probabilities, subject to source currency checks. \\
Horizon and source difficulty & Longer horizons have higher Brier (\(\rho=+0.161\), 95\% CI \([+0.026,+0.290]\)); Polymarket is harder than Manifold, which is harder than yfinance for every family. & Design clue; not enough alone for an applied rule. \\
YES underprediction & On contracts resolving YES (\(N=26/42\)), every family's mean \(p_{\text{success}}\) is below \(0.5\). & Supports calibration diagnostics. \\
Panel agreement and naive ensembles & Cross-family forecast correlations are high (mean \(\rho=+0.72\)); mean/median-of-five have higher Brier than the best single family on the \(N=142\) paired set. & Warns against naive averaging. \\
\bottomrule
\end{tabular}
\caption{Cross-family regularities that motivated calibration and allocation tests. They are useful only when converted into a scored rule and compared with simple baselines.}
\label{tab:applied-patterns}
\end{table}

The composed four-rule recipe converted these patterns into a scored forecast by applying a discount for very small probabilities, a middle-band YES correction, a horizon shrinkage, and a source-difficulty shrinkage to the mean panel forecast. The recipe beats naive aggregation on the same \(N=142\) paired set, but it does not clear the stronger best-single and evaluation bars balanced by source.

\begin{table}[htbp]
\centering
\small
\begin{tabular}{l c c c c}
\toprule
strategy & Brier & $\Delta$ vs.\ best-single & 90\% CI & $p_{\mathrm{perm}}$ \\
\midrule
median-of-5             & $0.272$ & $+0.018$ & -- & -- \\
mean-of-5               & $0.261$ & $+0.007$ & -- & -- \\
best-single (Claude)    & $0.254$ & --      & -- & -- \\
adjusted aggregate (universal-only) & $0.232$ & $-0.022$ & $[-0.050, +0.005]$ & $0.18$ \\
adjusted aggregate (per-channel)    & $0.246$ & $-0.009$ & $[-0.035, +0.016]$ & $0.58$ \\
\bottomrule
\end{tabular}
\caption{Composed adjustment on the \(N=142\) paired public-domain set. The within-cohort win is against naive aggregation, not against the strongest simple correction.}
\label{tab:composed-adjustment}
\end{table}

Three points matter. First, the adjusted aggregate improves over naive aggregation within this cohort: it beats median-of-five by \(-0.040\) Brier (\(p=0.0013\)) and mean-of-five by \(-0.029\) (\(p=0.0069\)). Second, the comparison with the best single family remains underpowered: the point estimate is favorable (\(-0.022\)) but the paired test is \(p=0.18\), with an estimated \(N\approx250\)--\(300\) required for the observed effect. Third, later audits balanced by source weaken the composite. The current source+\(\sigma\) allocation rule and diagnostic-triggered allocation have higher Brier than simpler rules; Hedge over raw-family, calibrated-family, and simple-pool experts is directionally better but nonsignificant (\(0.226481\) versus \(0.233529\), \(p=0.4671\)) and regresses on Manifold in the balanced slice. Choosing the best family in hindsight remains far better (\(0.117454\) Brier), so there is headroom across model families, but the observable rules have not recovered it.

The pairwise ranking experiment supports a controlled relative-judgment use. A four-family contrastive comparison supports pairwise ranking over 24 unique non-tie pairs (accuracy \(0.750\), utility \(+0.583\), \(p=0.0044\) versus random). Translation tests are favorable but not yet ready as a probability layer: translated probabilities beat the calibrated baseline by \(-0.022714\) Brier within the experiment, with \(p=0.0628\); one transfer direction clears \(p=0.0314\), the reverse direction misses at \(p=0.0636\); and the joined market control slice is only 24 rows with translated-vs-market \(p=0.5783\). A prospective Polymarket comparison has frozen 24 pairwise market bars before model calls, but unresolved outcomes cannot support a current scoring claim. Pairwise ranking therefore appears in the paper as a controlled use, while the current probability rule remains the calibration rule for very small probabilities.

\textbf{The correction for very small probabilities as a standalone rule beats raw on every family at $p<0.05$.} The rule applies only when the panel forecast is below \(0.20\), shrinking that low forecast toward \(0.10\) by \(\hat p_{\mathrm{low}} = 0.35\bar p + 0.65(0.10)\). The simpler form of the recipe, the single discount for very small probabilities applied per-family without the other three rules, improves per-family Brier on every panel member at $p < 0.05$:

\begin{table}[htbp]
\centering
\small
\begin{tabular}{l c c c c}
\toprule
family & raw Brier & discounted Brier & $\Delta$ & paired $p_{\mathrm{perm}}$ \\
\midrule
claude        & $0.2543$ & $\mathbf{0.2240}$ & $-0.0302$ & $0.016$ \\
codex-$5.5$   & $0.2625$ & $0.2416$         & $-0.0208$ & $0.030$ \\
codex-$5.4$-mini & $0.2714$ & $0.2450$       & $-0.0264$ & $0.015$ \\
deepseek      & $0.3222$ & $\mathbf{0.2704}$ & $\mathbf{-0.0518}$ & $\mathbf{0.0008}$ \\
gemini        & $0.3167$ & $0.2840$         & $-0.0327$ & $0.008$ \\
\bottomrule
\end{tabular}
\caption{Per-family correction for very small probabilities. The single rule improves Brier for each tested family on the public-domain panel, but later source currency checks restrict its use to eligible forward-looking rows.}
\label{tab:low-prob-family}
\end{table}

Discounted Claude at Brier \(0.2240\) is at parity or better than the four-rule adjusted aggregate (\(0.232\)), so the additional features add little on this corpus. The single adjustment for very small probabilities improves Brier at \(p<0.05\) across every model family, including DeepSeek and Gemini, which were not the panel that originated the rule. A fitted-calibrator audit using the stored rows did not identify a better replacement: source-isotonic slightly improved the overall point estimate but lost Manifold and was nonsignificant, while tail-beta shrinkage was worse. The 2026-06-04 source-documented rerun excludes 10 yfinance/yfinance\_etf complete panels lacking label-time documentation, leaving 132 panels; source-isotonic remains unsupported at \(-0.005248\) Brier versus the correction for very small probabilities with paired \(p=0.7099\), while raw mean-panel is \(+0.029598\) worse. A later source currency stress audit narrows the rule: it improves post-cutoff rows (\(-0.025326\), tail-only \(-0.101306\)) and regresses pre-cutoff/source-visible rows (\(+0.035016\), \(p=0.0002\); tail-only \(+0.097719\), \(p=0.0002\)). A row-level rerun through the shared source currency check leaves those scores unchanged while exposing 39/240 stored-flag-vs-computed-relation conflicts. The rule for very small probabilities is therefore forward-looking calibration tied to computed cutoff and label-time documentation, not retrospective benchmark correction.

\section{Controlled use under source and market constraints}
\label{sec:controlled-use}

Equal-information market controls set the scoring problem. Once the row, label, and comparator information state are fixed, the question becomes whether a named model-derived output improves a scored task under the same controls. Three current uses pass that bar within the paper's scope: the correction for very small probabilities improves forward-looking rows that pass the source currency check and fails on source-visible rows; pairwise comparisons rank contracts better than chance but are not yet a probability model; and one Gemini expert-training prompt beats bare and placebo prompts in a balanced public question comparison. A fourth line, family choice, has measurable headroom but no reliable selection rule. These results belong in the same paper because each depends on the same row validity checks that make the market controls interpretable.

The controlled-use map separates current uses from design evidence and negative guidance. Table~\ref{tab:applied-outputs} states what can be used now, what only motivates the next experiment, and what current controls leave unsupported.

\begin{table}[htbp]
\centering
\begingroup
\setlength{\tabcolsep}{3pt}
\renewcommand{\arraystretch}{0.95}
\small
\begin{tabular}{@{}L{0.23\linewidth} L{0.34\linewidth} L{0.33\linewidth}@{}}
\toprule
Output or signal & Current role & Boundary \\
\midrule
Calibration for very small probabilities & Deterministic post-processing rule for forward-looking rows with very low panel probabilities that pass the source currency check. & Regresses on source-visible rows; not a retrospective benchmark correction. \\
Pairwise ranking & Ranking or tournament interface when the task is relative difficulty or relative likelihood. & Probability translation and market-additive use still need larger and prospective controls. \\
Expert-training prompt & One Gemini public question intervention that beats bare, placebo, and same row calibrated bare forecasts. & Does not beat the market on current overlap; the 591/600-call Claude run remains underpowered and below gate, and Codex+DeepSeek does not reproduce it. \\
Structured fields & Candidate interface for forcing comparable evidence fields instead of relying on rationale prose. & Mixed small tests; design evidence only until a larger placebo-controlled packet clears. \\
Headroom across model families & Evidence that different families contain conditional signal worth selection or review. & Current cheap selection rules do not recover the hindsight headroom. \\
Negative controls & Evidence that generic reflection, self-revision, and selective action are weak prompt-only interventions. & Does not rule out tool-using, retrieval-grounded, expert-written, or heldout-tuned systems. \\
\bottomrule
\end{tabular}
\endgroup
\caption{Controlled-use map. The paper's applied contribution is a bounded account of which model signals are usable now, which remain design evidence, and where the present evidence says to stop.}
\label{tab:applied-outputs}
\end{table}

\paragraph{A scored-use procedure.} The current evidence supports a simple order of operations for any new forecast packet. First, attach forecast time, model cutoff, source, label vintage, event family key, and comparator timestamp to each row; rows that fail source currency or label-time checks can be studied diagnostically, but are excluded from broad score comparisons. Second, score any equal-information market or human baseline before reporting model gains. When that baseline wins, the model result is reported as diagnostic unless a predeclared calibrated rule or blend beats it on the same rows. Third, on forward-looking rows that pass the source currency check, apply the correction for very small probabilities only to very low panel probabilities and report the uncorrected score beside it. Fourth, use pairwise comparisons for prioritization or ranking; absolute-probability translation waits for prospective checks against raw, calibrated, and market baselines. Fifth, treat prompt variants as candidates only when they beat bare, placebo, calibrated bare, and source-split checks. This procedure is deliberately narrow, but it turns the market-negative results into a usable rule for when to score, adjust, rank, or leave model output alone.

\textbf{The rule for very small probabilities is the current selected probability rule.} The discount for very small probabilities is the best practical rule in the database, but only under conditions where the source currency check passes. It was selected after earlier calibration work, so the paper treats it as a scoped rule requiring independent replication rather than as a general principle. It improves every tested family on the public-domain panel, including families that were not used to originate the rule. On source-documented rows the raw mean-panel remains worse than the correction for very small probabilities by \(+0.029598\) Brier with paired \(p=0.0062\), BH \(q=0.0163\), and BY \(q=0.0616\). On the source currency stress panel, the same rule improves post-cutoff rows by \(-0.025326\) while regressing pre-cutoff/source-visible rows by \(+0.035016\) with \(p=0.0002\), BH \(q=0.0024\), and BY \(q=0.0091\). This is the pattern in miniature: the rule helps on the forward-looking slice and damages retrospective/source-visible rows. It is therefore a calibration result for rows that pass the source currency check, not a universal correction.

\textbf{Pairwise ranking is the current relative-judgment use case.} The strongest positive ranking evidence supports relative judgment rather than absolute probability. A comparison of same source minimal pairs, balanced by source, gives 24 unique non-tie pairs with accuracy \(0.750\), utility \(+0.583\), \(p=0.0044\) versus random, and \(p=0.0002\) versus source control. Later probability-translation tests are favorable in one direction, while the required single-contract probability checks remain unmet: within-experiment translated-vs-calibrated and translated-vs-raw checks, bidirectional transfer, joined market control, and prospective causal-order resolution are not all satisfied. The supported use is pairwise ranking or tournament support.

\textbf{Expert-training prompting is a Gemini-specific candidate.} The completed public question prompt comparison tests five conditions on the same 120 contracts across FRED, Manifold, and Polymarket: bare prompt, length-matched placebo, expert-training prompt, audit-informed prompt, and failure-mode-specific prompt. In the Gemini run, expert-training improves paired Brier versus bare prompt by \(-0.060569\) (63 wins, 29 losses, 28 ties; sign \(p=0.0005\), BH \(q=0.0031\), BY \(q=0.0115\)) and versus length-matched placebo by \(-0.024287\) (60 wins, 33 losses, 27 ties; sign \(p=0.0067\), BH \(q=0.0163\), BY \(q=0.0616\)). Mean Brier improves in all three sources. The two other structured prompt variants do not beat placebo. An external-control audit using stored forecasts adds two boundaries. First, expert-training still beats the calibrated bare prompt on the same rows by \(-0.052503\) Brier (64 wins, 33 losses, 23 ties; sign \(p=0.0022\), BH \(q=0.0103\), BY \(q=0.0390\)). Second, it does not beat markets on the current overlap: expert-training minus all matched market rows is \(+0.150950\) Brier over 51 matched rows (\(p=0.0110\), BH \(q=0.0239\), BY \(q=0.0904\)), and expert-training minus equal-information market rows is \(+0.093130\) over 33 matched rows (\(p=0.0351\), BH \(q=0.0601\), BY \(q=0.2271\)). A 591/600-call Claude run has 112 complete contract blocks: expert-training is directionally better than the bare prompt on mean Brier by \(-0.003638\) (40 wins, 44 losses, 31 ties; sign \(p=0.7436\), BH \(q=0.7436\), BY \(q=1.0000\)) and directionally better than length-matched placebo by \(-0.004175\) Brier (53 wins, 34 losses, 25 ties; sign \(p=0.0530\), BH \(q=0.0848\), BY \(q=0.3202\)). It remains underpowered and below the replication gate: it does not pass either sign test or the source split. Manifold is directionally favorable, while FRED regresses versus bare and Polymarket regresses versus placebo. The audit-informed Claude arm has favorable mean deltas against bare and placebo (\(-0.006726\) and \(-0.005266\)), but neither sign test is significant and the source split remains fragile. A staged Codex+DeepSeek check now has 448 scored calls (89 complete five-condition family-contract blocks; 90 expert-training paired blocks: 41 Codex and 49 DeepSeek). Across both families, expert-training is worse than bare prompt on mean Brier by \(+0.007173\), does not pass the sign test (\(p=0.4704\), BH \(q=0.5942\), BY \(q=1.0000\)), is worse than placebo by \(+0.065223\) (\(p=0.3891\), BH \(q=0.5187\), BY \(q=1.0000\)), and fails the source split because only FRED is directionally better than placebo. Codex remains directionally favorable on mean Brier in the current slice but is not clean across sources; DeepSeek regresses. This leaves the completed Gemini finding as a narrow candidate: one model family, scored after outcomes were known, weaker than matched market rows, with the Claude run underpowered and below gate and Codex+DeepSeek not reproducing the effect.

\textbf{Structured fields and family choice remain design evidence.} Some small tests favor typed evidence fields over free prose, but the placebo-control continuation is negative for the stronger two-step claim. Family choice has real headroom: choosing the best family in hindsight reaches much lower Brier than current decision rules, and family-by-contract interaction is substantial. Current observable selection rules do not recover that headroom: Hedge over raw, calibrated, and simple-pool experts is directionally better but not significant, and graph-family weighting is small and fragile. These results identify where model information appears; they do not yet supply a reliable rule for choosing among model families.

\textbf{Generic prompt intervention is still weak evidence.} Generic reflective prompting, selective action, and self-revision do not reliably improve forecasts under controls. The tested generic variants either fail confirmation, overcorrect, produce no measurable change, or underperform calibration for very small probabilities. The Gemini expert-training comparison shows that a prompt can improve Brier under a tighter design, but it does not rescue the broader claim that asking the model to reason harder, repair itself, or act selectively is enough to improve forecast probabilities.

These results are analyzed together because the positive findings depend on the same validity checks that make the market controls interpretable. The common unit is a forecast row whose information state is documented before any calibration, ranking, or selection rule is evaluated, not an isolated prompt or model family.

\section{Secondary diagnostics: bias transfer and prompt stability}
\label{sec:bias-transfer}
\label{sec:prompt-invariance}
\label{sec:freq-framing-test}

The bias-transfer and prompt-stability experiments are not the paper's central contribution. They appear in the main text as diagnostics for why naive prompt interventions and family-general allocation do not yet support decisions.

\textbf{Bias transfer.} Three novel-bias tests on the same \(N=30\) public-domain contracts across five families found a structured split: loss-frame and probability-weighting probes were close to symmetric or near-linear, while current-state/status-quo framing moved several families substantially. A held-out seven-bias slate later supported the representational distinction between text-discussed motivational cases and direct utility-like mechanisms. The anti-bias-prompt follow-up was weakened by raw-gap controls, and the out-of-distribution slate was inconclusive and confounded by cross-family pretraining differences. We therefore use these results only to interpret family/source sensitivity and prompt-risk, not to claim that cognitive-debiasing prompts improve forecasting skill.

\textbf{Prompt stability.} The prompt-invariance audits give the same warning from another angle. Standard forecast prompts yield relatively stable point estimates across all five families, but reference-class outside-view prompting is a different reasoning architecture rather than a harmless paraphrase. The worry channel is much noisier across prompts than \(p_{\text{success}}\), and cross-prompt worry-Brier signs can reverse relative to the low-overlap experiment. Any decision rule that depends on a verbal uncertainty channel must therefore be treated as family-, corpus-, and prompt-bound until it clears heldout controls.

These diagnostics bound the controlled-use result: model output does not become reliable through generic reflection, debiasing, or self-revision. It can be used only when the interface is tied to a scored channel, a calibration on rows that pass the source currency check rule, a pairwise ranking task, or a predeclared allocation rule that beats simpler controls.

\section{The re-audit discipline: a warning to the field}
\label{sec:reaudit}

The discipline that produced the results above also produced a re-audit of the study's own prior findings. We treat the re-audit itself as a contribution because the same overcalling pattern is a plausible risk in LLM-forecasting studies with small per-source or per-family cells.

\subsection{Power calculus applied to our own study}

Before applying the power-aware audit, eight prior ``no effect'' or ``cross-corpus replication'' claims in the study had been written down as findings. Applying the Fisher-$z$ audit retrospectively:

\begin{itemize}[leftmargin=*]
\item Four ``null'' findings move from ruled out to underpowered: the data did not actively rule out a meaningful effect at the per-family target $|\rho|$ floor; the apparent lack of effect was mainly a lack of power.
\item Five cross-corpus claims move from ``replicated'' or ``failed-to-replicate'' to corpus-specific: the original cross-corpus reading was a single sub-source point estimate at $N{=}5$--$15$ external, with CIs that crossed zero and with sub-source heterogeneity unmodeled.
\item One self-knowledge-inversion finding (originally three families at $N{=}5$) retracts to a single family at $N{=}15$ on the low-overlap corpus; the other two cross-corpus directions flip on rerun.
\item Two per-family decision rules (Brier-$\Delta$ map and \emph{consider-both-sides} effect) move from usable to ``directional hypothesis only''; the required $N$ to confirm a $\Delta_{\mathrm{Brier}}{=}0.05$ effect at 80\% power is $\approx 459$ per family, an order of magnitude above the original $N{\approx}30$ per cell.
\end{itemize}

In aggregate, $\approx 50\%$ of the study's prior conclusions changed under the power-aware re-audit. None of the underlying executions were flawed. The unsupported step was the earlier interpretation.

\subsection{A worked example: the per-family worry sign-flip retest}

The previous Gemini worry-channel hypothesis was scope-restricted to ``low-overlap corpus, original prompt only'' pending a $\sim$42-call retest of the same prompt on the public-domain $N{=}42$ slice used by the other public-domain tests. The retest, run with the same worry-only prompt on \texttt{gemini-2.5-flash} via API:

\begin{itemize}[leftmargin=*]
\item Schema-OK: $42/42$.
\item $\rho(\mathrm{worry}, \mathrm{Brier}) = +0.110$ -- the direction-correct direction, not the inverted direction the prior rule had assumed.
\item Fisher-$z$ detectable $|\rho|$ at $N{=}42$, $80\%$ power: $\geq 0.43$. Observed $|\rho|{=}0.11$ is well below this; the CI excludes the original ``inverted'' direction ($\rho < -0.30$).
\item The inversion hypothesis is ruled out at the target effect size; the direction-correct point estimate remains underpowered for the positive hypothesis.
\end{itemize}

The Gemini channel rule stays withdrawn. The earlier ``inversion'' claim, estimated from $N{=}5$ with $N_{\mathrm{wrong}}{=}2$, was a sample-size effect. This is the lowest-cost decisive retest available; two comparable retests remain open.

\subsection{Implications for the published LLM-forecasting literature}

Two prominent $2024$ contributions to LLM forecast calibration are (a)~the ensemble-prediction result that a 12-LLM aggregate matches human-crowd Brier on real binary questions~\cite{schoenegger2024ensemble} and (b)~the retrieval-augmented LLM that approaches human-superforecaster accuracy on a contamination-resistant slice~\cite{halawi2024approaching}. Both report claims at sample sizes that, under the same Fisher-$z$ power calculus we apply to our own study, are limited by:

\begin{itemize}[leftmargin=*]
\item For a Brier-delta $\Delta{=}0.05$ comparison (LLM ensemble vs.\ human crowd) at $80\%$ power, paired permutation requires $N{\geq}300$--$500$ per side.
\item For a per-model rank correlation $|\rho|{=}0.30$ at the same power: $N{\geq}91$.
\item For a per-(model, sub-source) cell at $|\rho|{=}0.40$: $N{\geq}50$ per cell -- requiring sub-source $N{\geq}250$ across five families when sub-sources are five-way mixed.
\end{itemize}

We do not re-analyse those datasets here; we observe only that the standard ``2024-style'' aggregate claim is structurally compatible with the same overcalling failure mode our own study exhibited. Paleka et al.~\cite{paleka2025pitfalls} make the broader version of the same warning: LLM forecasting evaluation is unusually vulnerable to temporal leakage, unreliable date-restricted retrieval, benchmark extrapolation errors, and circular comparisons against human forecasts already visible to the model or retrieval system. Consistency-check benchmarks offer a complementary way to probe coherence before outcomes resolve~\cite{paleka2024consistency}, but they do not remove the need for resolved, equal-information scoring. The methodological gap this paper targets is a consistent power-aware audit: every claim labelled with one of three possible outcomes, $n_{\mathrm{required}}$ pre-computed before launch, sub-source heterogeneity reported alongside the aggregate, and prior verdicts re-audited under the same calculus as new ones.

Recent work has expanded the comparison set since those papers. Table~\ref{tab:lit-positioning} summarizes the recent benchmark classes and the boundary this paper draws around them.

\begin{table}[p]
\centering
\begingroup
\setlength{\tabcolsep}{2pt}
\renewcommand{\arraystretch}{0.86}
\scriptsize
\begin{tabular}{@{}L{0.20\linewidth} L{0.28\linewidth} L{0.21\linewidth} L{0.23\linewidth}@{}}
\toprule
Class & Examples & What they measure & Boundary in this paper \\
\midrule
Future-question and system benchmarks & ForecastBench, Prophet Arena, and AIA Forecaster~\cite{karger2024forecastbench,yang2025prophetarena,alur2025aiaforecaster} & Forecast generation, expert comparison, and system performance on live or future questions. & The comparison row still needs source currency, label-time, and equal-information documentation. \\
Belief updating & EVOLVECAST~\cite{yuan2025evolvecast} & Whether forecasts move appropriately when new post-cutoff information is supplied. & The information state at the original forecast time is part of the row definition, not an afterthought. \\
Numerical forecast intervals & QuantSightBench~\cite{qin2026quantsightbench} & Prediction intervals for continuous future quantities, scored by coverage and interval sharpness. & Interval calibration is a different output interface; the present paper studies binary-event rows, equal-information baselines, and controlled use after row validation. \\
Question generation and resolution & Automated forecasting-question generation and resolution~\cite{bosse2026automatingforecasting} & Generation and resolution of forecasting questions for AI evaluation. & Automatic resolution still needs label-time, settlement-rule, and equal-information-comparator metadata before row scores support model-vs-comparator claims. \\
Confidence elicitation and fictional markets & Confidence elicitation and fictional prediction-market framing~\cite{xiong2024confidence,todasco2025fakepredictionmarkets} & Whether models can emit useful confidence or wager-like signals about answer correctness. & A confidence or stake signal becomes forecasting evidence only after it is joined to resolved rows and equal-information baselines. \\
Trading and replay benchmarks & Prediction Arena, PolyBench, and PredictionMarketBench~\cite{zhang2026predictionarena,cheng2026polybench,arora2026predictionmarketbench} & Profit, loss, execution, fees, liquidity, timing, and position sizing. & Trading profit is not the same object as same-contract Brier under equal information. \\
Market-style evaluation and coordination & Foresight Arena, MarketBench, and Reppo~\cite{nechepurenko2026foresight,fradkin2026marketbench,reppo2026} & Market-style scoring, coordination, or infrastructure for AI forecasting and training data. & The market comparison still needs a dated baseline and enough resolved rows to distinguish small edges. \\
Market relationship discovery & Semantic Trading~\cite{capponi2025semantictrading} & Agentic discovery of correlated, anti-correlated, and otherwise linked prediction-market pairs. & Market-pair structure can motivate relative judgment, but probability use still needs resolved outcomes and equal-information controls. \\
Relative judgment & Strategic Foresight venture tournament~\cite{csaszar2026strategicforesight} & Pairwise model rankings against human managers and investors. & Supports pairwise ranking as a credible evaluation format, but not standalone probability translation. \\
Evaluation warnings & Pitfalls and consistency checks~\cite{paleka2025pitfalls,paleka2024consistency} & Temporal leakage, extrapolation risk, circular comparisons, and pre-resolution coherence. & This paper adds scored database audits and equal-information market controls. \\
\bottomrule
\end{tabular}
\endgroup
\caption{Related benchmark classes and the boundary used in this paper. The present study is not a live-trading or autonomous-system benchmark; it audits the row-level evidence needed before model, human, and market comparisons are interpreted as forecasting evidence.}
\label{tab:lit-positioning}
\end{table}

Read against that literature, the present paper is not a live-market benchmark or an autonomous-trading evaluation. Its distinct contribution is the row-level documentation needed to interpret those benchmarks: before comparing an LLM, a market, and a human, the row must specify what was source-visible, what label vintage is admissible, and whether the baseline was measured under equal information. This is also why evidence that markets score better is scientifically useful. If an LLM has higher Brier after the information state is equalized, the remaining question is whether any model-derived outputs remain useful under calibration, ranking, and family/source constraints.


\paragraph{What an audit would measure.} The evidence in this paper is enough to show that row-level validity can change conclusions in this program. It is not yet enough to claim that the same failure rate holds across the field. An audit would treat each benchmark row as the unit of analysis and record the following checks before comparing model, human, or market scores:

\begin{table}[p]
\centering
\begingroup
\scriptsize
\setlength{\tabcolsep}{2pt}
\renewcommand{\arraystretch}{0.9}
\begin{tabular}{L{0.24\linewidth} L{0.43\linewidth} L{0.25\linewidth}}
\toprule
Audit check & What must be recorded & Why it matters \\
\midrule
Source currency & Forecast timestamp, model cutoff or retrieval window, resolution date, and whether the resolved answer was source-visible at generation time. & Separates future-event prediction from retrieval or source familiarity. \\
Label-time validity & Outcome label, label source, data vintage, and any later revisions or settlement-rule changes. & Prevents scoring against values that were not admissible at resolution time. \\
Equal-information comparator & Human, crowd, market, or agent baseline measured on the same contract at the same pre-outcome information time. & Prevents comparing a model forecast with a comparator that knew more or less. \\
Effective sample size & Event-family identifiers, source strata, repeated sibling markets, and family/model repetition. & Prevents row-rich but event-thin conclusions. \\
Decision rule status & Whether the tested rule was predeclared, tuned on the same rows, or evaluated prospectively. & Separates exploratory diagnostics from supported use. \\
\bottomrule
\end{tabular}
\caption{Broader validity audit protocol. The present manuscript supplies this protocol, a scored within-program audit, and the public-audit status summarized in Table~\ref{tab:public-benchmark-audit}; a cross-benchmark failure-rate claim requires applying the protocol to several public benchmark families.}
\label{tab:field-audit-protocol}
\endgroup
\end{table}

For reproducibility, \path{projects/llm_forecasting_calibration_program/tools/field_wide_validity_audit_protocol.py} materializes this protocol as a row schema and seed matrix over twelve benchmark or evaluation families. A source inventory classifies the current external access surface: ForecastBench and PredictionMarketBench are high-access row-level routes, four routes are medium-access, and six still require a public trace or row-release check before scoring. Table~\ref{tab:public-benchmark-audit} summarizes the public routes already checked in this paper.

\begin{table}[htbp]
\centering
\scriptsize
\begin{tabular}{@{}L{0.18\linewidth} L{0.21\linewidth} L{0.34\linewidth} L{0.19\linewidth}@{}}
\toprule
Route & Row-level status & Completed check & Remaining limit \\
\midrule
ForecastBench & Public forecast rows and question rows available. & 500 question rows inspected; 475 have core validity fields; 250 have timestamped same-contract market rows. The 2026 processed-forecast audit scores 70 files over 521 unique resolved binary row keys and 230 event family keys; 68 files have an equal-information market slice. Only 6 files beat the prior-day market baseline before and after event family capping; the median file-level delta is \(+0.0866\) Brier points in both views. A 2024 human-comparator audit scores 141 files over 7{,}259 row keys and 766 event family keys; the human-super/public aggregate Briers are \(0.1186/0.1532\). & Official-data rows still need data-vintage documentation; the human aggregate files each have only two strict equal-information market rows, so this is not a broad human/market comparison. \\
Prediction\hspace{0pt}Market\hspace{0pt}Bench & Public replay episodes available. & Four replay episodes contain 33 settled tickers, 378{,}596 orderbook rows, 297{,}273 trade rows, and 370{,}254 reconstructable same-time market-baseline rows. & Released episodes contain no stored LLM forecast rows; model comparisons require submitted agent traces or a new benchmark run. \\
Prophet Arena & Public sample task rows available. & Four public sample releases contain 68 task rows with task ids, source/event tickers, prediction deadlines, context, metadata close times, and 26 resolved rows; five public AI Prophet repositories were checked for trace files. & The fetched samples contain no submitted forecast probabilities, model input timestamps, or same-time market/human baseline probabilities, and no public Prophet Arena submission/leaderboard trace archive was found, so no conclusion-change score is available. \\
PolyBench & Repository and database schema available; released database unavailable in this run. & The source-access check verifies the public repository/schema surface, confirms zero GitHub releases and zero committed database/CSV/parquet row files, and records that the OneDrive link resolves to HTML rather than a direct database file. & A released database or equivalent row export is required before scoring. \\
Halawi 2024 binary-resolved benchmark & Local date summary only. & The local summary records the corpus-validity warning that no 2025-or-later resolutions are present in the stored date distribution. & The raw benchmark rows are not present locally, so this is not a completed external row audit. \\
\bottomrule
\end{tabular}
\caption{Current public benchmark audit status. These checks strengthen the measurement-validity route, but they are not evidence of prevalence across the field.}
\label{tab:public-benchmark-audit}
\end{table}

The public-audit results are deliberately limited. ForecastBench supports public forecast-score and human-comparator audits, PredictionMarketBench supports a public replay-row market-baseline audit, and Prophet Arena supports task-row source-access checks. The ForecastBench human aggregate files are useful because they are scoreable and public, but their equal-information market overlap is only two rows per aggregate file. Prophet Arena and PredictionMarketBench still need submitted forecast traces before model-vs-baseline conclusions can be recomputed; the public Prophet Arena repositories checked here expose task rows and evaluation code, not a submission archive. PolyBench and Halawi currently remain data-access cases. These files therefore support the row-level validity argument and specify the missing joins; they do not yet establish a cross-benchmark failure rate.

\subsection{Source-currency, label-time, and equal-information audit}

The central re-audit result is that benchmark validity is a moving target. A dataset that was a fair future-event benchmark for one model generation can become a retrieval or source-familiarity benchmark for a later generation. The local claim register records a Halawi 2024 binary-resolved date histogram with no 2025-or-later resolutions; for current model generations, the corresponding eligibility filter, \emph{resolve\_date $>$ max(panel\_cutoff)}, would leave zero rows in that local summary. This is not a replacement for a raw-row external audit, but it is the concrete corpus-validity drift warning that motivates the broader validity protocol.

\begin{table}[htbp]
\centering
\begingroup
\scriptsize
\setlength{\tabcolsep}{2pt}
\renewcommand{\arraystretch}{0.9}
\begin{tabular}{L{0.24\linewidth} L{0.43\linewidth} L{0.25\linewidth}}
\toprule
Check & Current evidence & Conclusion \\
\midrule
Manifold source currency panel & 80 matched contracts / 240 tool-free calls; post-minus-pre Brier $+0.191098$, paired-stratum delta $+0.2155$, permutation $p=0.0004$, BH $q=0.0031$, BY $q=0.0115$. Partial market-prior repair joins 51 contracts and leaves the effect positive under adversarial missing-band sensitivity. & Manifold measurement result: source-visible and post-cutoff rows cannot be pooled without documentation; source-general prevalence remains unmeasured. \\
Manifold pre-outcome market bar & On the 51 repaired contracts, market Brier is $0.099673$ versus joined LLM-panel $0.166963$ overall; market+LLM blending is unsupported ($p=0.794$). These rows are typed as not-equal-information external baselines. & Useful diagnostic comparison, but not a broad human/crowd baseline. \\
Polymarket source check & Gemini and DeepSeek aggregate post-minus-pre directions are positive, but matched source/topic/length strata are null or opposite-sign. The initial seven-day market-price freeze design only fills 4/24 rows; horizon sweep tops out at 12/24. & Diagnostic evidence, not replication across sources or a usable equal-information design. \\
Polymarket replacement equal-information sample & Replacement sampling on 2026-06-15 selects 24 one-per-event rows from 80 eligible candidates at a two-day freeze horizon, 14 NO / 10 YES, all open-by-target with nonempty CLOB history. The four-family panel has higher Brier than the market: panel $0.267758$ vs market $0.072964$, paired $p=0.0068$, BH $q=0.0163$, BY $q=0.0616$. & Lower market Brier in this slice under the raw paired test and BH correction; BY treats it as sensitivity rather than arbitrary-dependence significance at \(0.05\). \\
Manifold equal-information sample & Public API history fill validates 24/24 rows, 15 NO / 9 YES. Selected five-family panel Brier is $0.198723$ vs Manifold $0.160977$, panel-minus-market $+0.037746$, paired $p=0.5431$, BH $q=0.6207$, BY $q=1.0000$. & Second-source comparison; the market Brier is lower, but the paired test is inconclusive. \\
Manifold same-day freeze expansion & Public API history fill validates 32/34 request rows, with two unsupported or unfetched rows excluded, 15 NO / 17 YES. Selected five-family panel Brier is $0.214665$ vs Manifold $0.135951$, panel-minus-market $+0.078714$, paired $p=0.0048$, BH $q=0.0163$, BY $q=0.0616$. & Larger same-contract Manifold check; market Brier is lower under the raw paired test and BH correction, but BY treats it as sensitivity and this is still not a population estimate. \\
Manifold horizon sensitivity & Public API history fills validate 18/18 one-day, 10/10 two-day, and 7/7 seven-day request rows. Selected five-family panel versus Manifold market Brier is $0.202270$ vs $0.099699$ at one day ($p=0.0122$, BH $q=0.0244$, BY $q=0.0921$), $0.231846$ vs $0.109365$ at two days ($p=0.0152$, BH $q=0.0281$, BY $q=0.1060$), and $0.228263$ vs $0.193649$ at seven days ($p=0.5045$, BH $q=0.6054$, BY $q=1.0000$). & Overlapping sensitivity check only; these rows are not used as independent core market-control evidence. \\
FRED official-data check & Fixed one-year FRED series supplies 49 pre-cutoff rows; full paired post-minus-pre delta is only $+0.016477$ ($p=0.30375$). In the vintage repair, 15 of 98 binary labels changed. Across 192 blinded-control calls, the apparent current-label post-minus-pre penalty collapses from $+0.024719$ to $-0.002989$ after vintage repair. & Official-data diagnostic; label-time checks are required before current-label positives count. \\
Metaculus target cells & Current authenticated endpoints expose post/question payloads but not resolved binary values plus dated aggregate history for the sampled rows; data-download endpoint is restricted. & Remains a data-access question, not a negative result. \\
\bottomrule
\end{tabular}
\caption{Re-audit evidence by source. The paper's broad comparison claims depend on source currency, label-time validity, and equal-information status, not merely on model call volume.}
\label{tab:reaudit-sources}
\endgroup
\end{table}

The source currency panel is matched on source, topic, question-length bucket, and cutoff relation. Reliable base-rate bands are not available for that panel, so base-rate matching remains a stated limitation rather than a hidden control. The database now contains 170 external market baseline rows: 51 not-equal-information Manifold diagnostic rows, 4 equal-information Polymarket rows from the failed seven-day design, 24 equal-information Polymarket rows from the replacement sample, 24 equal-information Manifold rows from the 2026-06-15 history fill, 32 equal-information Manifold rows from the same-day freeze expansion, and 35 additional Manifold horizon-sensitivity rows at one, two, and seven days before resolution. The 51-contract Manifold market-prior join is useful as a diagnostic repair, but not as the paper's market comparison: market mean Brier is \(0.099673\), leave-one-out tuned market+LLM blend Brier is \(0.097218\), paired delta is \(-0.002455\) with \(p=0.794\), and the post-cutoff subset prefers market-only. Its effective unit is 51 contracts with outcome mix 17 YES / 34 NO, and 32 of 51 market rows have Brier below \(0.05\). The baseline availability audit also leaves 29 of the 80 source currency-panel contracts without a joined market or human comparison. The three primary equal-information slices are therefore the relevant market controls for this manuscript, with the smaller horizon fills used only as overlapping sensitivity checks. They do not support raw LLM superiority: Polymarket has much lower market Brier than the model panel, the first Manifold fill has lower market Brier but an inconclusive paired test, and the 32-contract Manifold expansion is again lower-Brier for the market under the raw paired test and BH correction, with BY sensitivity reported in Table~\ref{tab:multiplicity-audit}. The one-day, two-day, and seven-day Manifold horizon checks are reported only as sensitivity rows. Because these controls are still small and post-hoc, they are not used as a high-power filter for every downstream calibration or prompt claim; they are used to state that the current evidence does not support raw model panel superiority and to define the larger equal-information acquisition that the benchmark still needs.

This audit changes what the paper can support. The Manifold source-currency result is strong enough to report as a measurement-validity contribution for that source. The FRED check gives the clearest label-time example: scoring the same official-data question against a current value and against the admissible vintage can reverse the apparent pre/post difficulty contrast. The equal-information market rows do not support a broad reading in which LLMs are superior to markets or humans, and narrow the applied result: LLM-derived information has to be tested through source-specific calibration, pairwise ranking, structured interfaces, or future allocation rules that beat the market or human bar under the same information rule.

\subsection{Claim rule after re-audit}

The re-audit leaves a simple rule for interpreting applied results. A per-family or per-source decision rule is treated as supported only when the scored rows pass the source currency and label-time screens, the effective denominator is stated at the contract, pair, or event family level, and the result survives the relevant equal-information, placebo, source split, or held-out family control. Results that fail one of those checks remain useful as diagnostics or entries in the continuation matrix, but they do not become present-tense applied claims. This is why the manuscript keeps the correction for very small probabilities, pairwise ranking, and Gemini prompt result under explicit scope while leaving family selection, structured fields, source currency beyond Manifold, and broader market/human comparisons as next tests.

\section{What this paper does not establish}
\label{sec:limits}

\begin{itemize}[leftmargin=*]
\item It does not show that LLMs beat humans, human crowds, or prediction markets. The database has 170 external market baseline rows, including 119 equal-information market rows across Polymarket and Manifold. On the completed 24-contract Claude+Codex+Gemini+DeepSeek Polymarket replacement slice, the market Brier is lower than the four-family model panel under the raw paired test and BH correction, with BY sensitivity above \(0.05\). On the separate 24-contract Manifold fill, the market Brier is lower than the selected five-family low-stake model panel, but the paired comparison is inconclusive. On the 32-contract Manifold same-day freeze expansion, the market Brier is again lower than the selected five-family panel under the raw paired test and BH correction, with BY sensitivity above \(0.05\). The smaller Manifold one-day, two-day, and seven-day horizon checks are overlapping sensitivity rows, not additional independent market-control evidence. Broad superiority would require predeclared or sufficiently powered baselines balanced by source under the same pre-outcome information rule.
\item It does not estimate a general market control effect from the three equal-information slices. The slices are strict but still modest. They do not support a raw superiority claim in this dataset, especially on Polymarket and the 32-contract Manifold expansion, but they are not enough to characterize market and model performance across sources, horizons, liquidity regimes, or event families.
\item It does not establish translated pairwise probabilities as a standalone probability layer. Pairwise ranking survives as source heldout evidence with promising translation tests, but within-experiment, transfer, market control, and prospective causal-order checks remain incomplete.
\item It does not establish that the source-currency result generalizes beyond the main Manifold panel. Polymarket tests are aggregate-positive but matched-stratum null/opposite-sign, and FRED current-label positives weaken under label-time/vintage repair.
\item It does not establish a prompt method that generalizes across sources or beats markets. The completed Gemini public question comparison supports one expert-training prompt against bare, placebo, and calibrated bare forecasts on the same rows, but generic reflection, selective action, and self-revision fail or regress in other audits. Current matched market rows remain stronger. The 591/600-call Claude run remains underpowered and below the replication gate: expert-training is directionally better than bare and placebo on mean Brier, but does not clear the sign tests and does not pass the source split. The audit-informed arm is only directionally favorable and fragile across sources. A 448-call Codex+DeepSeek staged replication is worse than bare and placebo on mean Brier and does not clear the sign-test or source-split checks. The prompt result is therefore a Gemini-specific candidate until a second completed family run clears the same checks.
\item It does not establish external generality or provider-independent generality. The low-overlap corpus is private and not publicly available in raw form; the low-overlap frequency-framing, bid-ask, worry-sign, and channel findings need public replication at proper power. The current low-overlap/public split also confounds novelty, source, topic, question length, and implicit base rates, so a four-cell de-confounded corpus remains the main methodological follow-up. The scored calls use proprietary APIs or CLIs; replication on open models and public questions is required before claiming provider-independent generality.
\item It does not establish a working rule for choosing among model families. Best-family-in-hindsight scores show that different families make different errors, but choosing the best family after seeing outcomes is not an intervention. Headroom across model families becomes useful only if observable features, outside review, market disagreement, or held-out model signals recover that headroom prospectively after cost.
\item It does not decompose LLM reasoning mechanistically or identify the causal role of post-training. Channel-orthogonality results are expression-level measurements, not white-box interpretability. Family identity predicts several bias and channel differences, but the comparisons are cross-family and confound post-training with pretraining data. Isolating an alignment effect requires within-family checkpoints or matched-pretraining families.
\end{itemize}

\paragraph{Effective denominators.} The main results are not counted at the model call level. Repeated model calls over one contract improve a panel estimate, but they do not create new events; this prevents model call rows from being treated as independent events. Table~\ref{tab:effective-denominators} records the denominator used for the central slices; the support audit also records source counts and evidence paths.

\begin{table}[p]
\centering
\begingroup
\setlength{\tabcolsep}{2.5pt}
\renewcommand{\arraystretch}{0.9}
\scriptsize
\begin{tabular}{@{}L{0.18\linewidth} L{0.18\linewidth} L{0.19\linewidth} L{0.35\linewidth}@{}}
\toprule
Evidence slice & Scored rows & Denominator used & Boundary \\
\midrule
Source-currency panel & 240 calls over 80 contracts & Contract & Manifold-supported measurement result; extension beyond Manifold still needs non-Manifold matched pre/post rows. \\
Equal-information market controls & 24 Polymarket contracts, 24 Manifold contracts, a 32-contract Manifold same-day expansion, and smaller Manifold horizon sensitivity checks & Contract at a matched market time & Market controls for this manuscript, not a broad market/human result; horizon checks are overlapping sensitivity rows. \\
Prompt comparison & 600 calls over 120 contract-condition blocks & Contract-condition block & Gemini-specific public question result and replication target; current market overlap has lower market Brier; the 591/600-call Claude run remains underpowered and below gate, and Codex+DeepSeek does not replicate it. \\
Pairwise ranking & 94 calls over 24 non-tie pairs & Unique pair & Ranking result only; probability translation and prospective market freeze scoring remain incomplete. \\
FRED label-time repair & 196 calls over 98 series/event rows & FRED series/event row & Label-time diagnostic; vintage repair changes labels but does not create an improvement rule. \\
ForecastBench public score audits & 2026 score audit: 70 forecast files, 521 row keys / 230 event family keys. 2024 human-comparator audit: 141 files, 7{,}259 row keys / 766 event family keys; human aggregate files have 577 resolved rows each. & Public row key, with event family-capped market check & Public-audit feasibility and human-comparator coverage; human aggregate market overlap is only two rows per file, so this is not a broad human/market comparison. \\
\bottomrule
\end{tabular}
\endgroup
\caption{Effective denominators for the central evidence slices.}
\label{tab:effective-denominators}
\end{table}

The full support audit listed in Table~\ref{tab:reproduction-status} gives the per-slice source and event-group status.

\paragraph{Next tests.} Table~\ref{tab:next-tests} states the evidence that would change each claim's status. The most direct route to a broader measurement contribution is an audit of public forecasting benchmarks and market-replay environments. That route does not require LLMs to beat markets; it requires showing that row-level source currency, label-time, or equal-information checks are often missing and can change conclusions. A broad LLM-vs-market or LLM-vs-human claim requires a predeclared or substantially larger equal-information sample that beats the market or human baseline under the same pre-outcome information rule, not more calls on the current rows.

The immediate unblocking protocol for the \(N=24\) market control bottleneck has two admissible paths. The first is strict backfill: mine existing pre-outcome model calls only when the same contract has an auditable market-history export at or before the model timestamp or a predeclared freeze timestamp; keep one row per event family unless a clustered analysis is reported; and reject rows with unresolved YES-token mapping, final-page odds, or post-outcome price selection. A 100-row non-Polymarket export pass using the current local database surfaced only 10 eligible Manifold request rows, so local backfill alone does not solve the denominator problem. The second path is prospective acquisition: freeze market prices before model calls, hide the prices from the prompts, persist the prompt/output packet before resolution, and score only after outcomes settle. The target for a stronger market control result is \(N \ge 100\) resolved contracts, balanced by source across at least two market sources, with event family clustered uncertainty and predeclared comparisons among market-only, raw panel, calibration on rows that pass the source currency check, pairwise-derived rankings or probabilities where applicable, and market+model blends.

The immediate experimental roadmap has two jobs. First, turn the current post-hoc equal-information controls into a predeclared \(N \ge 100\) packet, using strict backfill where timestamps are auditable and prospective freezes where they are not. The new Manifold horizon fills show that local acquisition can add valid rows, but also that the remaining denominator is source- and event-family-limited. Second, test whether the Gemini expert-training result generalizes by completing a full second-family packet. The current staged evidence makes Codex-only completion more informative than more mixed Codex+DeepSeek aggregation; an open model is the next independent replication target. A positive prompt result is useful only if it beats bare, placebo, calibrated bare, and same-time market or human baselines on the same rows and does not vanish in the source split.

A source currency beyond Manifold claim requires Metaculus/export access or another non-Manifold panel with matched pre/post resolution-date coverage and label-time documentation. Pairwise ranking becomes a standalone probability layer only if within-experiment, transfer, market control, and prospective causal-order checks all clear. The prompt result becomes broader only if a structured-prompt arm replicates in another completed model-family run and beats a larger market or human overlap measured at the same time. The low-overlap corpus findings become externally general only after a sanitized release or a public niche-domain replication that breaks the current novelty/source/topic/length/base-rate confound. Provider-independent generality requires repeating the public question validity and controlled-use checks on open models. A separate evidence matrix records each candidate result, the checks it has passed, the checks still missing, and the next decisive test. A follow-up priority matrix ranks the follow-up tests by claim impact, minimum next step, what would strengthen each claim, and what would rule it out.

\begin{table}[p]
\centering
\begingroup
\setlength{\tabcolsep}{2.5pt}
\renewcommand{\arraystretch}{0.86}
\tiny
\begin{tabular}{@{}L{0.18\linewidth} L{0.22\linewidth} L{0.30\linewidth} L{0.20\linewidth}@{}}
\toprule
Claim & Present treatment & Evidence that would change the status & If the evidence is absent or negative \\
\midrule
Benchmark validity beyond this study & Protocol, within-program audit, ForecastBench public score and human-comparator audits, Prophet Arena source-access pilot, PredictionMarketBench row-schema pilot, and PolyBench source-access pilot. & Row-level audit of several public benchmark families, recording source currency, label-time validity, equal-information status, and conclusion changes after repair. & Keep claims about prevalence across the field out of the paper. \\
LLM vs.\ market/human performance & Not supported; current equal-information slices have lower market Brier or are inconclusive, and the strict controls are still small and partly overlapping. & Predeclared, equal-information sample with at least \(N \ge 100\) resolved rows, event family clustered uncertainty, and market/history timestamps frozen before model scoring. & Keep raw-panel superiority claims out of the paper. \\
Source currency beyond Manifold & Supported on the main Manifold panel; other sources remain diagnostic. & Non-Manifold panel with matched pre/post rows, admissible label vintage, and base-rate documentation. & Treat source currency as a Manifold-supported measurement result plus a general audit requirement. \\
Pairwise ranking as probability layer & Ranking use supported; probability translation remains provisional. & Within experiment and transfer replication plus prospective probability translation against raw, calibrated, and market controls. & Use pairwise comparisons only for ranking or tournament support. \\
Structured prompting intervention & Expert-training improves paired Brier versus bare, length-matched placebo, and calibrated bare forecasts on the same Gemini rows; two other prompt variants do not beat placebo; current market overlap has lower market Brier; the 591/600-call Claude run remains underpowered and below the replication gate, while Codex+DeepSeek does not reproduce it. & Complete a full second-family packet, prioritizing Codex-only completion over more mixed Codex+DeepSeek aggregation; then run an open-model replication; require survival by source and larger market or human overlap measured at the same time. & Treat the current result as a Gemini-specific candidate interface. \\
Structured evidence fields & Hypothesis only. & Larger balanced paired test that beats free prose, same-turn fields, and two-call prose after runtime failures are included. & Treat structured fields as a future paired-test design until that comparison is run. \\
Family allocation & Best-family-in-hindsight headroom exists; current allocation rules do not recover it. & Predeclared observable features, outside review, market disagreement, or held-out model signals that recover best-family headroom after cost. & Report headroom as a design target without claiming a working selector. \\
External generality of low-overlap results & Not yet established. & Sanitized release or public niche-domain replication that breaks the current novelty/source/topic/length/base-rate confound. & Keep low-overlap results as secondary diagnostics. \\
Open model replication & Not yet established. & Repeat the public question source currency, market control, calibration for very small probabilities, and pairwise-ranking checks with open models. & Treat current model-family results as provider-snapshot evidence. \\
\bottomrule
\end{tabular}
\endgroup
\caption{Claim-status table. The paper separates current claims from the evidence that would change their status; missing or negative follow-up evidence leaves the manuscript at its present measurement-validity and controlled-use scope.}
\label{tab:next-tests}
\end{table}

\section{Conclusion}
\label{sec:conclusion}

The main lesson is practical: scored probabilities do not become forecasting evidence until the row says what the model could have known. Each row must specify what was source-visible to the model generation, whether the label uses an admissible time vintage, and whether the human or market comparison was measured under the same information rule. In this study, those checks change the claim: the Manifold source currency panel supports a measurement-validity result, while the equal-information Polymarket and Manifold market comparisons do not support a raw model panel superiority reading of the current evidence and do not estimate a general effect comparing markets and models.

The applied claim is limited but nontrivial because it identifies where model information survives strong controls. The current evidence supports calibration for very small probabilities and pairwise ranking as a relative-judgment task. It also records one Gemini expert-training prompt result, structured evidence fields, and headroom across model families as follow-up targets rather than general deployed methods. Generic reflection, self-revision, and selective action are not enough. The resulting order of use is concrete: screen rows before scoring, apply the correction for very small probabilities only on eligible forward-looking rows, use pairwise comparisons for relative triage rather than standalone probabilities, treat the expert-training prompt as a Gemini replication target, and leave generic reflection, self-revision, selective-action prompts, and model-family selection unsupported until they clear the same controls.

A broader claim requires new evidence rather than more calls on the same rows. The decisive tests are a broader row validity audit, prospective or larger equal-information samples, non-Manifold source currency panels with label-time documentation, and predeclared methods that beat raw, calibrated, and market controls under the same information state. The market control upgrade is especially concrete: strict backfill is allowed only with auditable pre-outcome market histories, and the cleaner route is a prospective freeze packet with at least 100 resolved rows, source balance, event family clustering, hidden market prices, and predeclared scoring. These requirements also define the companion benchmark design: row validity metadata, equal-information baselines, calibrated probabilities, relative judgments, prompt interventions, choosing among model families, open model replication, and public low-overlap substitutes are scored as separate tracks rather than collapsed into one aggregate leaderboard. The same design is useful during research, not only after publication, because it tells a live packet what comparison it is allowed to answer before the outcome is scored.

\section{Reproducibility}
\label{sec:repro}

All model calls persist as JSONL files and are ingested into the SQLite database at \path{analytics/public/calibration/forecaster_calibration.db}. Pre-registrations and verdict resolutions are in the same database. The reusable general-purpose statistics module is at \path{src/ztare/experiment_stats.py} (power calculator, bootstrap CI, paired permutation, Fisher-$z$ Spearman, TOST, BH-FDR, BY robustness, power-aware verdict, BIC Bayes factor, reproducibility hash). Forecasting-specific wrappers are in \path{src/ztare/forecasting/calibration_stats.py}; the public methodology file lists the exact commands for each finding.

\begin{flushleft}
\footnotesize
Project scripts directory: \path{projects/llm_forecasting_calibration_program/tools/}\\
Figure scripts directory: \path{llm-forecast-calibration-cross-corpus/evidence/reproducers/}
\end{flushleft}

\textbf{Reproducing the scope checks.} Three audit scripts define the scope boundary:

\begin{itemize}[leftmargin=*]
\item \texttt{\detokenize{paper_readiness_exhaustion_audit.py}}
\item \texttt{\detokenize{paper_coherence_audit.py}}
\item \texttt{\detokenize{independent_equal_information_source_audit.py}}
\end{itemize}

Equal-information acquisition is split into deterministic scripts:

\begin{itemize}[leftmargin=*]
\item \texttt{\detokenize{equal_information_baseline_export_packet.py}}
\item \texttt{\detokenize{equal_information_baseline_result_ingest.py}}
\item \texttt{\detokenize{equal_information_freeze_feasibility_audit.py}}
\item \texttt{\detokenize{equal_information_horizon_sweep.py}}
\item \texttt{\detokenize{equal_information_replacement_sample_acquire.py}}
\item \texttt{\detokenize{equal_information_replacement_dispatch_packet.py}}
\item \texttt{\detokenize{equal_information_replacement_score.py}}
\item \texttt{\detokenize{non_polymarket_equal_information_export_packet.py}:} emits Manifold equal-information request rows.
\item \texttt{\detokenize{non_polymarket_equal_information_result_acquire.py}:} fills Manifold request rows from public market history.
\item \texttt{\detokenize{non_polymarket_equal_information_result_ingest.py}:} ingests the validated Manifold rows.
\item \texttt{\detokenize{non_polymarket_equal_information_score.py}:} scores the joined Manifold model-vs-market comparison.
\item \texttt{\detokenize{field_wide_validity_audit_protocol.py}:} emits the row schema and 12-route seed matrix for the broader validity-audit route.
\item \texttt{\detokenize{field_wide_validity_source_inventory.py}:} records the current external-source access status for those 12 routes and names the next row-level extraction step for each.
\item \texttt{\detokenize{field_wide_forecastbench_row_schema_pilot.py}:} applies the row schema to the local ForecastBench 2026-04-12 question bundle and reports validity-field coverage.
\item \texttt{\detokenize{field_wide_forecastbench_score_audit.py}:} scores public ForecastBench processed-forecast files and compares eligible market-source rows with the prior-day market value; the manuscript uses it for the 2026 score audit and the 2024 human-comparator audit.
\item \texttt{\detokenize{field_wide_prophet_arena_row_schema_pilot.py}:} fetches public Prophet Arena sample releases, audits their task rows against the paper's row schema, and checks public AI Prophet repositories for submitted forecast trace files.
\item \texttt{\detokenize{field_wide_predictionmarketbench_row_schema_pilot.py}:} inspects public PredictionMarketBench replay episodes and reconstructs same-time market-baseline rows from orderbook snapshots and settlements.
\item \texttt{\detokenize{field_wide_polybench_source_pilot.py}:} verifies the public PolyBench repository/schema surface, GitHub release/file status, and noninteractive OneDrive dataset status.
\item \texttt{\detokenize{field_wide_validity_local_evidence.py}:}\\
Emits the provenance-limited Halawi date-distribution summary used as a warning, not as a completed audit.
\item \texttt{\detokenize{claim_gap_matrix.py}:} emits an evidence matrix separating supported results from underpowered results, claims not valid for broad conclusions, and claims requiring external data.
\item \texttt{\detokenize{decisive_continuation_matrix.py}:} ranks the follow-up tests by claim impact, minimum next step, what would strengthen the claim, and what would rule it out.
\item \texttt{\detokenize{evidence_upgrade_plan.py}:}\\
Emits a clean evidence-upgrade plan separating the public-benchmark audit route, the structured-prompting test, and the larger equal-information baseline sample.
\item \texttt{\detokenize{experiment_coverage_summary.py}:} reads the historical findings ledger and current queue, then emits a reader-facing count summary for included, compressed, deferred, and excluded experiment families.
\item \texttt{\detokenize{applied_signal_coverage_audit.py}:} records each applied signal as supported, bounded, diagnostic, negative guidance, or a future evidence route, with a manuscript anchor and next check.
\item \texttt{\detokenize{scored_use_procedure_audit.py}:} checks that the scored-use procedure has manuscript text, support files, and a stop condition for each operational step.
\item \texttt{\detokenize{prospective_counterexplanation_design_audit.py}:} checks that each planned positive result has a counter-explanation and a before-scoring design check before it can support a broader claim.
\item \texttt{\detokenize{reviewer_concern_coverage_audit.py}:} checks that the likely reviewer concerns have explicit manuscript answers, evidence files, remaining boundaries, and paper anchors.
\item \texttt{\detokenize{forecast_row_validity_benchmark_blueprint.py}}:\\
Converts the claim-gap and continuation matrices into a companion public-benchmark design, including row validity fields, equal-information comparators, applied tracks, and failure conditions.
\item {\footnotesize\path{llm-forecast-calibration-cross-corpus/evidence/benchmark/run_benchmark.py}}:\\
Validates and scores a small row-contract packet under the paper's validity rules. The validator requires model-family and evaluation-track fields, separates source-visible rows, label-time failures, eligible forecast rows, and equal-information comparator rows, then reports track and family counts, calibration deltas, pairwise accuracy, and one-row-per-event-family comparator summaries. It is a companion tool, not a new result.
\item \texttt{\detokenize{numeric_claim_trace_audit.py}:} checks the manuscript's headline numerical statements against the SQLite database and stored score reports.
\item \texttt{\detokenize{central_evidence_effective_n_audit.py}:} records calls, contracts or pairs, market rows, source counts, and event-group documentation for the central evidence slices.
\item \texttt{\detokenize{multiple_testing_effective_n_audit.py}:} collects the p-values used in the manuscript into one BH-FDR family, adds BY robustness values for arbitrary dependence, and records the effective denominator for each test.
\item \texttt{\detokenize{literature_positioning_audit.py}:} checks that each related-work class in Table~\ref{tab:lit-positioning} has a bibliography key, source URL, and explicit paper boundary.
\item \texttt{\detokenize{submission_readiness_audit.py}:} checks required sections, table/figure labels, generated support files, evidence-count floors, broad-claim boundaries, and LaTeX log health.
\item \texttt{\detokenize{rendered_pdf_smoke_audit.py}:} extracts the compiled PDF text and checks freshness, page count, required rendered sections, boundary sentences, and absence of internal draft language.
\item \texttt{\detokenize{make_equal_information_figure.py}:} regenerates Figure~\ref{fig:equal-info-bars} from the stored equal-information market control score reports.
\end{itemize}

Together these tools let a reader verify the current market-baseline coverage, inspect the replacement sample before forecasts, reproduce the first same-contract market comparison, score the public ForecastBench processed-forecast round under the paper's row checks, inspect Prophet Arena public task-row access and submitted-trace status, inspect PredictionMarketBench replay-row validity, inspect the PolyBench source-access gap, trace the paper's headline numbers to current score files, audit central denominators and event-group coverage, recompute the consolidated multiplicity table, audit related-work positioning, audit how the historical experiment log was compressed, inspect which applied signals are supported or bounded, check the scored-use procedure and its stop conditions, check the before-scoring counter-explanations for planned positive results, check where likely reviewer concerns are answered, inspect the companion benchmark blueprint implied by the missing evidence, run a minimal row-validity benchmark packet end to end, verify that the compiled PDF still contains the required claim boundaries, rank the decisive follow-up tests, inspect the evidence-upgrade plan, and see exactly what evidence is still missing for a broad human/crowd claim, broader validity claim, or stronger intervention claim.

\textbf{Reproduction status.} Table~\ref{tab:reproduction-status} separates what is reproducible now from what requires a sanitized or substitute release. The private low-overlap corpus is the main unreleased component, but it affects only the secondary low-overlap elicitation findings. The source currency, label-time, equal-information, market control, and calibration claims on rows that pass the source currency check are represented by public-market, official-data, database, scoring, and audit machinery in the repository.

\begin{table}[p]
\centering
\begingroup
\setlength{\tabcolsep}{2.5pt}
\renewcommand{\arraystretch}{0.86}
\tiny
\setlength{\abovecaptionskip}{3pt}
\setlength{\belowcaptionskip}{0pt}
\begin{tabular}{L{0.27\linewidth} L{0.31\linewidth} L{0.34\linewidth}}
\toprule
Component & Current status & Reproduction role \\
\midrule
SQLite evidence database & Present at \path{analytics/public/calibration/forecaster_calibration.db}. & Reproduces call counts, score joins, source currency screens, 165 label-time rows, and market-baseline coverage. \\
Scoring/audit scripts & Present in the project tool directories and listed above. & Recomputes the paper's readiness, coherence, equal-information, and label-time checks without new model calls. \\
Evidence, follow-up, and upgrade matrices & Present; generated by the claim-gap, continuation, and evidence-upgrade scripts listed above. & Lists candidate results, missing checks, interpretation changes, and next actions. \\
Applied signal coverage audit & Present; generated by the applied-signal script listed above. & Lists applied components, evidence, use case, boundary, next check, and manuscript anchor. \\
Scored-use procedure audit & Present; generated by the scored-use script listed above. & Checks the operational steps for row screening, equal-information baselines, scoped calibration, ranking-only pairwise use, and prompt-variant gates. \\
Prospective counter-explanation design audit & Present; generated by the prospective counter-explanation script listed above. & Checks that planned positive results name the simpler explanation they must rule out before scoring. \\
Reviewer concern coverage audit & Present; generated by the reviewer-concern script listed above. & Maps likely reviewer concerns to the manuscript answer, evidence files, remaining boundary, and paper anchor. \\
Forecast row validity benchmark blueprint and runner & Present as both a design report and a runnable row-contract validator under \path{llm-forecast-calibration-cross-corpus/evidence/benchmark/}. & Specifies and tests the companion benchmark shape implied by the missing evidence: row validity, equal-information comparators, calibration, relative judgment, intervention, choosing among model families, open model replication, and public low-overlap substitute tracks. The bundled example exercises explicit model-family and track fields, calibration output, a pairwise judgment, source-visible rows, and a label-time failure. \\
Broader audit protocol and pilots & Present; generated by the broader-audit scripts listed above. & Supplies the row schema, 12-route source inventory, ForecastBench audits, Prophet Arena task-row check, PredictionMarketBench replay-row coverage, PolyBench source-access status, and Halawi warning; it is not evidence of prevalence across the field. \\
Numeric claim trace & Present; generated by the numeric-trace script listed above. & Checks headline numerical statements against the current database and stored score reports. \\
Central evidence denominator audit & Present; generated by the effective-N script listed above. & Records calls, contracts or pairs, market rows, source counts, and event-group status for the main evidence slices. \\
Multiplicity and effective-denominator audit & Present; generated by the multiple-testing script listed above. & Reports all p-values used by the manuscript in one BH-FDR family, adds BY robustness values, and records the denominator for each test. \\
Literature positioning audit & Present; generated by the literature-positioning script listed above. & Checks that related-work classes have bibliography keys, source URLs, and a clear boundary against the present paper's claim. \\
Submission-readiness audit & Present; generated by the submission-readiness script listed above. & Checks manuscript structure, generated support files, evidence-count floors, broad-claim boundaries, and build-log health. \\
Rendered-PDF smoke audit & Present; generated by the rendered-PDF script listed above. & Checks the compiled PDF, not only the source, for freshness, required rendered sections, boundary text, and internal-language leaks. \\
Public-market samples & Present for Polymarket and Manifold equal-information comparisons. & Reproduces the market control boundary claims. \\
Open-model replication & Not yet run. & Needed for provider-independent generality; not required for the paper's current provider-snapshot claims. \\
Raw low-overlap questions & Not publicly releasable in raw form. & Needed only for direct replication of the private low-overlap elicitation findings. \\
Sanitized or substitute low-overlap corpus & Planned release path: neutralized identifiers plus a public niche-domain substitute with the same four-axis profile. & Enables third-party replication of the frequency-framing, bid-ask, and low-overlap channel findings without exposing private workflow details. \\
\bottomrule
\end{tabular}
\endgroup
\caption{Reproduction status. The central source currency and equal-information claims are reproducible from public project files; the low-overlap channel findings require a sanitized or substitute corpus.}
\label{tab:reproduction-status}
\end{table}

\bibliographystyle{plain}
\bibliography{refs}

\appendix

\section{Evidence ledger for compressed diagnostics}
\label{app:evidence-ledger}

The main text compresses several experiment families for readability. Table~\ref{tab:evidence-ledger} records what was compressed and what insight is retained. The detailed run records remain in the SQLite database, the project workspace, and the public methodology/claim-summary files referenced in $\S\ref{sec:repro}$.

\begin{table}[htbp]
\centering
\small
\begin{tabular}{L{0.23\linewidth} L{0.35\linewidth} L{0.34\linewidth}}
\toprule
Compressed family & Retained insight & Where it is used \\
\midrule
Uncertainty channels & Worry, spread, trajectory variance, self-predicted Brier, and outside-view base rates expose distinct side information, but their sign and value are family/source conditional. & Motivates source/family conditioning and prevents a universal channel-rule claim. \\
Self-assessed channel choice & Several families misidentify which emitted uncertainty channel predicts their own error. & Supports the negative self-monitoring result and the need for scored interfaces. \\
Universal calibration regularities & Overconfidence on very small probabilities, horizon/source difficulty, YES underprediction, and high family agreement recur across families. & Supplies the candidate features tested by the rule for very small probabilities and composed adjustment. \\
Composed adjustment and model-family headroom & The adjusted aggregate beats mean/median in-cohort but fails stronger scrutiny balanced by source; choosing the best family in hindsight remains much better than observed allocation rules. & Separates real headroom from established allocation evidence. \\
Pairwise ranking and translation & Pairwise contrastive ranking is stronger than raw probability translation; translation remains promising but is not yet supported as a standalone probability model. & Retains pairwise ranking as a controlled use, not evidence that LLMs are superior to markets or humans. \\
Source and label-time audits & Manifold source currency, FRED vintage repair, failed Polymarket freeze design, replacement Polymarket, and Manifold equal-information fill show that validity checks change conclusions. & Forms the measurement-validity contribution and bounds the applied claims. \\
Prompt intervention & Expert-training improves paired Brier versus bare and placebo prompts in a completed Gemini public question comparison; generic reflection, self-revision, and diagnostic allocation mostly fail or produce no measurable change under controls. & Keeps a Gemini-specific candidate for replication while rejecting broad prompt-only improvement. \\
\bottomrule
\end{tabular}
\caption{Evidence-preservation table for compressed diagnostics. Compression changes placement, not the underlying scope boundaries.}
\label{tab:evidence-ledger}
\end{table}

\section{Coverage audit for omitted or deferred work}
\label{app:coverage-audit}

The project log contains many more experiments than the main paper reports. The inclusion rule is conservative: a result appears in the main text only if it changes the validity layer, a supported use case, a stated limit, or a continuation test. The \path{projects/llm_forecasting_calibration_program/tools/experiment_coverage_summary.py} script parses both the full research log and the curated findings ledger before compression. In the current snapshot it detects 111 unique numbered rows in the full research log, with 74 outside the curated paper ledger. The curated ledger then classifies 37 paper-relevant rows: 3 central validity/control/calibration rows, 19 secondary diagnostics retained in the paper, 7 retractions/supersessions/underpowered boundaries, 7 sibling workflow or non-forecasting findings excluded from this manuscript, and 1 execution/persistence row with no paper claim. The \path{projects/llm_forecasting_calibration_program/tools/applied_signal_coverage_audit.py} script separately records 8 applied components: 6 supported or scoped candidate components, 1 negative guidance row, and 1 external evidence route. The structured-prompt comparison is now complete at 600/600 scored Gemini calls across 120 contracts; the expert-training arm passes the bare/placebo and calibrated bare checks on the same rows, while market checks remain unfavorable, the 591/600-call Claude run remains underpowered and below the replication gate, and the 448-call Codex+DeepSeek staged run does not pass the replication gate. Table~\ref{tab:coverage-audit} records the main excluded or deferred families so that compression does not hide evidence.

\begin{table}[htbp]
\centering
\small
\begin{tabular}{L{0.27\linewidth} L{0.33\linewidth} L{0.32\linewidth}}
\toprule
Work family & Why it is not a main claim & How the insight is preserved \\
\midrule
Prompt-intervention variants & One expert-training prompt beats bare and placebo prompts in the completed Gemini public question comparison; generic reflection, selective action, self-revision, and diagnostic-triggered allocation do not beat the relevant controls. & Reported as a Gemini-specific candidate plus negative evidence against unvalidated prompt-only interventions. \\
Objective effort / coding-task calibration & This is a sibling problem about effort prediction and hidden-test performance, not the same forecast row validity question. & Excluded from this paper unless it later supplies a forecasting-specific intervention that clears source and market controls. \\
Proof-audit and workflow-only findings & These improve the research workflow but ask reviewers to switch domains away from event forecasting. & Excluded from the main manuscript; relevant only as provenance for the audit discipline. \\
Low-overlap elicitation retests & Several channel findings are promising but still corpus-bound or underpowered for claims across sources. & Preserved as diagnostics and continuation tests: replicate on a public niche corpus or sanitized release before generalizing. \\
Fitted calibrators and allocation rules & Source-isotonic, graph-family weighting, and diagnostic allocation have not beaten simpler controls robustly. & Reported as headroom evidence; applied use waits for panels balanced by source or an external reviewer/market expert. \\
Prospective market freeze comparisons & Frozen market bars exist for some future comparisons, but unresolved outcomes cannot score current claims. & Listed as continuation tests; no standalone probability layer until outcomes resolve and market, raw, and calibrated controls are passed. \\
Deconfounded corpus design & The current low-overlap/public split confounds novelty, source, topic, and horizon. & Treated as the key methodological follow-up for external generality, not as evidence already in hand. \\
\bottomrule
\end{tabular}
\caption{Coverage audit for work compressed out of the main text. The table separates omission for irrelevance, omission for weak evidence, and deferral because the right outcome data or control does not yet exist.}
\label{tab:coverage-audit}
\end{table}

\section{Composed adjustment recipe: closed form}
\label{app:composed-adjustment-recipe}

Let $\bar p$ be the mean of the five families' point estimates $p_{\text{success}}$ on a given contract, $h$ the days-to-resolution-from-cutoff, and $s\in\{\textrm{polymarket}, \textrm{manifold}, \textrm{yfinance}\}$ the source identifier. The composed adjusted forecast $\hat p$ is:

\begin{align*}
\hat p_{\mathrm{low}} &= \begin{cases} 0.35\,\bar p + 0.65\,(0.10) & \text{if } \bar p < 0.20 \\ \bar p & \text{otherwise} \end{cases} \\
\hat p_{\mathrm{YES\text{-}bias}} &= \hat p_{\mathrm{low}} + \begin{cases} +0.06 & \text{if } 0.30 \le \hat p_{\mathrm{low}} \le 0.55 \\ 0 & \text{otherwise} \end{cases} \\
\hat p_{\mathrm{horizon}} &= 0.5 + (\hat p_{\mathrm{YES\text{-}bias}} - 0.5)\cdot \tfrac{1}{1 + 0.01\,h} \\
\hat p &= 0.5 + (\hat p_{\mathrm{horizon}} - 0.5)\cdot w_s
\end{align*}

with $w_s = 0.70$ for polymarket, $0.85$ for manifold, $1.00$ for yfinance. The four coefficients $(0.10, 0.06, 0.01, w_s)$ are not learned from held-out folds; they are chosen from the regularities reported in $\S\ref{sec:universal}$ (overconfidence on very small probabilities, middle-band YES underprediction, horizon-conditional Brier slope, per-source Brier ordering). The per-channel alternative described in the main text substitutes per-family channel weights for the universal $w_s$; its weights are listed in the appendix support file. The implementation is in the forecasting-calibration workspace and reproduces on the database snapshot dated 2026-05-28.

\end{document}
